\documentclass[11pt,a4paper]{article}
\usepackage[margin=25.5mm,headheight=14pt]{geometry}
\usepackage[T1]{fontenc}
\usepackage{lmodern,amsmath,amssymb,amsthm,mathtools,mathrsfs,booktabs,longtable,enumitem,hyperref,graphicx,fancyhdr,xcolor,listings,setspace}
\usepackage[nopatch=footnote]{microtype}
\setstretch{1.045}
\hypersetup{pdftitle={Exact Hoeffding Constants for Finite Exchangeable Sequences},pdfauthor={Zinan Wang},colorlinks=true,linkcolor=blue!45!black,citecolor=blue!45!black,urlcolor=blue!40!black}
\pagestyle{fancy}\fancyhf{}\fancyhead[LO]{\small WANG -- SHARP FINITE EXCHANGEABILITY}\fancyfoot[C]{\thepage}
\newtheorem{theorem}{Theorem}[section]
\newtheorem{proposition}[theorem]{Proposition}
\newtheorem{lemma}[theorem]{Lemma}
\newtheorem{corollary}[theorem]{Corollary}
\newtheorem{remark}[theorem]{Remark}
\newtheorem{definition}[theorem]{Definition}
\newtheorem{conjecture}[theorem]{Conjecture}
\DeclareMathOperator{\Var}{Var}\DeclareMathOperator{\Cov}{Cov}\DeclareMathOperator{\Hyp}{Hyp}\DeclareMathOperator{\Bern}{Bern}\DeclareMathOperator{\supp}{supp}\DeclareMathOperator{\artanh}{artanh}
\newcommand{\E}{\mathbb E}\newcommand{\R}{\mathbb R}\newcommand{\N}{\mathbb N}\newcommand{\one}{\mathbf1}
\newcommand{\psiH}{\psi_{N,k,m}}
\newcommand{\CN}{C_N^\star}
\lstset{basicstyle=\ttfamily\scriptsize,breaklines=true,columns=fullflexible,frame=single,rulecolor=\color{black!20},keepspaces=true,showstringspaces=false}
\title{\bfseries Exact Hoeffding Constants for Finite Exchangeable Sequences}
\author{Zinan Wang\thanks{Copyright \textcopyright\ 2026 Zinan Wang. All rights reserved.}}
\date{11 October 2026}
\begin{document}
\maketitle
\begin{abstract}
We determine the least constant in the Hoeffding-type exponential-moment inequality for arbitrary centered linear contrasts of bounded exchangeable sequences of length $N$. The sharp finite-population coefficient is
\[
 C_N^\star=\begin{cases}
 \displaystyle N/(N-1),&N\text{ even},\\[3pt]
 \displaystyle[(N-1)\operatorname{artanh}(1/N)]^{-1},&N\text{ odd}.
 \end{cases}
\]
The even coefficient coincides with the maximal normalized variance. In odd populations the optimal Laplace parameter is nonzero, and the sharp coefficient strictly exceeds the variance benchmark. We obtain the result through a two-level reduction to hypergeometric laws, a global tilted-variance bound, and a mean-dependent quadratic envelope for their Laplace transforms. The envelope holds for both parities throughout the sub-half-population parameter range; parity enters when the remaining discrete margins are optimized. We also classify the odd-population extremizers and derive fourth-order expansions of the optimized hypergeometric proxies.
\end{abstract}
\noindent\textbf{Keywords.} Hypergeometric distribution; finite exchangeability; optimal sub-Gaussian proxy; Fisher noncentral hypergeometric law; Jacobi polynomials; exponential tilting; Riccati equation.\\
\textbf{MSC 2020.} 60E15, 60F10, 33C45, 62E15.
\section{Introduction}
\subsection{An exact finite-population concentration problem}
Let $X=(X_1,\ldots,X_N)\in[-1,1]^N$ be exchangeable and let $a\in\R^N$ satisfy $\sum_i a_i=0$. We seek the least constant $C_N^\star$ for which
\begin{equation}\label{eq:exchange-intro}
 \E e^{\lambda\langle a,X\rangle}
 \le\exp\left(\frac{\lambda^2C_N^\star}{2}\|a\|_2^2\right),
 \qquad \lambda\in\R,
\end{equation}
uniformly over the law of $X$ and the contrast $a$. Zero padding and orthogonal projection give the equivalent formulation for $\sum_{i=1}^n w_i(X_i-\bar X_N)$, where $\bar X_N=N^{-1}\sum_iX_i$. This formulation includes arbitrary signed weights.

For independent bounded variables the reference coefficient is one. Finite exchangeability introduces an exact population-dependent inflation. Barber \cite{Barber,BarberErratum} established weighted exchangeable Hoeffding bounds, and Lee and Kim \cite{LeeKim} obtained an upper coefficient $1+3/(2N)+O(N^{-2})$, a matching order-$N^{-1}$ lower bound, and an exact hypergeometric variational representation. We evaluate that representation for every $N$. The answer is $N/(N-1)$ when $N$ is even and $[(N-1)\operatorname{artanh}(1/N)]^{-1}$ when $N$ is odd.

Parity distinguishes two extremal mechanisms. The largest normalized variance is
\[
 D_N:=\sup_{X,a}\frac{\Var\langle a,X\rangle}{\|a\|_2^2}
 =\frac{4\lfloor N^2/4\rfloor}{N(N-1)}.
\]
For even populations, $C_N^\star=D_N$ and the sharp coefficient is determined at zero tilt. For odd populations,
\[
 D_N=\frac{N+1}{N}<C_N^\star<\frac{N}{N-1},
 \qquad C_N^\star-D_N=\frac{2}{3N^2}+O(N^{-3}).
\]
Thus the odd coefficient measures a genuine exponential-moment effect beyond the largest variance. A single coordinate centered by the population mean, $X_1-\bar X_N$, attains the odd bound at the explicit nonzero argument $\log((N+1)/(N-1))$ on a nearest-balanced Hamming slice. Proposition \ref{prop:variance_parity} and Proposition \ref{prop:explicit_extremizer} make this distinction precise.

\subsection{From exchangeable contrasts to finite-population sampling}
The central reduction identifies the exact finite-population problem with an optimization over hypergeometric random variables. The convex reduction to Hamming slices is followed by the two-level extremal principle of Lee and Kim, giving
\[
 \text{exchangeable centered contrasts}\ \longrightarrow\
 \text{hypergeometric optimal variance proxies}.
\]
We state and prove this variational identity immediately after the main results. In particular, a sharp estimate for each hypergeometric law can be transferred to every signed contrast of the original exchangeable sequence.

Two complementary estimates determine the optimum. First, a sharp bound on hypergeometric variance at every exponential tilt implies the universal coefficient $N/(N-1)$, which is attained for even populations. Second, a quadratic upper envelope for variance as a function of the tilted mean gives a scaled-Bernoulli Laplace comparison. The certificate for this envelope is valid in the sub-half-population regime for both parities. For odd $N$, all admissible success counts stay away from the exact half-population balance; optimizing the Bernoulli comparison then singles out the one-draw laws. For even $N$, the balanced zero-tilt equality accounts for the sharp constant. This separation of the continuous Laplace comparison from the discrete optimization isolates the role of parity.

\subsection{Extremal refinements and relation to prior work}
The exact constant leaves a finer question: how do nonsingleton samples approach the extremal value? For fixed $m$ and fixed odd success imbalance $s$, Theorem \ref{thm:odd_expansion} proves uniqueness of the global maximizing tilt and expands it through $N^{-3}$, while expanding the optimized proxy through $N^{-4}$. For the nearest-balanced population,
\[
 L_N-R_{N,1,m}
 =\frac{2(m-1)}{3N^3}
 +\frac{2(m-1)(2m+1)}{3N^4}+O_m(N^{-5}).
\]
The first two correction orders therefore agree with the singleton law; the order-$N^{-3}$ term resolves the competition between sample sizes. Theorem \ref{thm:uniform_fixed} additionally maximizes over every success count at fixed $m$, and Theorem \ref{thm:exact_imbalance} computes the exact quadratic response to a continuous success imbalance at each finite population size. The global localization of the maximizing tilt supplies the analytic remainders in these expansions.

Hoeffding \cite{Hoeffding}, Serfling \cite{Serfling}, Bardenet and Maillard \cite{BardenetMaillard}, and Greene and Wellner \cite{GreeneWellner} develop exponential inequalities and finite-sampling corrections. Diaconis and Freedman \cite{DiaconisFreedman} provide the classical finite-exchangeability framework. Gottschling and Caprio \cite{GottschlingCaprio} bound tails of sample means of bounded infinitely exchangeable sequences relative to the extremal means in the support of the de Finetti mixing measure. Equation \eqref{eq:exchange-intro} concerns finite sequences and arbitrary zero-sum contrasts. Fisher's noncentral hypergeometric family describes the exponential tilts used here \cite{Fog}.

Kou and Ying \cite{Kou} connect fixed-margin laws with Jacobi polynomials and sums of independent Bernoulli variables; Borcea, Br\"and\'en, and Liggett \cite{BBL} develop the broader stability framework. The Jacobi identities and zero-location theorem are classical \cite[Chapter IV and Section 6.7]{Szego}. Arbel, Marchal, and Nguyen \cite{ArbelMarchalNguyen} study optimal variance proxies and strict sub-Gaussianity, and Marchal and Arbel \cite{MarchalArbel} use differential equations for Beta and Dirichlet proxies. Atouani, Marchal, and Arbel \cite[Theorem 3.2 and Corollary 3.3]{AtouaniMarchalArbel} give the three-mass proxy formula specialized in Section \ref{sec:twodraw}. That section applies their formula to two draws and evaluates the resulting success-count extremum and singleton comparison. The all-degree variance certificates, the parity-uniform Laplace envelope, the classification of odd-population extremizers, and the optimized expansions for each fixed sample size are the hypergeometric results developed here. The fourth-order local moment expansion in \cite[Remark 3.6]{AtouaniMarchalArbel} is in the Laplace argument; our expansion is in $N^{-1}$ at the globally optimizing argument, and is accompanied by an extremum theorem over all success counts.

The main statements are followed immediately by the exchangeability-to-hypergeometric reduction. We then develop the curvature estimate, determine the parity-dependent sharp constants, and establish the unified Laplace envelope. The extremal expansions and local imbalance theory follow. Technical cumulant calculations, subsidiary curvature examples, analytic contraction details, and algebraic certificates are collected in the appendices.

\section{Statements of the main results}\label{sec:main}
For $N\ge2$ and $1\le k,m<N$, write
\begin{equation}\label{eq:hypergeom}
 \Pr\{H=i\}=\frac{\binom ki\binom{N-k}{m-i}}{\binom Nm},
 \qquad \max(0,m+k-N)\le i\le\min(k,m).
\end{equation}
Define the centered log-Laplace transform and its tilted variance by
\begin{equation}\label{eq:psi}
 \psi_{N,k,m}(t)=\log\E e^{t(H-km/N)},\qquad
 v_{N,k,m}(t)=\psi_{N,k,m}''(t).
\end{equation}
The second derivative is the variance of $H$ under the probability law having Radon--Nikodym derivative proportional to $e^{tH}$. The centering has no effect on this second derivative. Define $\sigma^2_{\mathrm{opt}}(H)=\sup_{t\ne0}2\psi_{N,k,m}(t)/t^2$. The quotient $2\psi_{N,k,m}(t)/t^2$ is assigned its continuous value $\Var(H)$ at $t=0$.

Denote by $\CN$ the least constant in \eqref{eq:exchange-intro} as the law of $X$, the centered vector $a$, and $\lambda$ all vary.
\begin{theorem}[Exact finite-exchangeability constants]\label{thm:exchange}
For every integer $N\ge2$,
\begin{equation}\label{eq:CNexact}
 \boxed{\quad C_N^\star=\begin{cases}
 \displaystyle\frac{N}{N-1},&N\text{ even},\\[5pt]
 \displaystyle\frac{1}{(N-1)\operatorname{artanh}(1/N)},&N\text{ odd}.
 \end{cases}\quad}
\end{equation}
In particular,
\begin{equation}\label{eq:CNupper}
 C_N^\star\le\frac{N}{N-1},
\end{equation}
with strict inequality for odd $N$. For odd populations, the maximum in the hypergeometric variational identity \eqref{eq:LKidentity} is attained precisely at
\begin{equation}\label{eq:odd_extremizers}
 m\in\{1,N-1\},\qquad k\in\{(N-1)/2,(N+1)/2\}.
\end{equation}
For both parities,
\begin{equation}\label{eq:asymptotic}
 \lim_{N\to\infty}N(C_N^\star-1)=1.
\end{equation}
\end{theorem}
\begin{proof}
The universal upper bound and even-population equality are established in Sections \ref{sec:exchangeable} and \ref{sec:parity}. The odd formula and equality cases follow from Theorem \ref{odd:thm:odd}; the limit is proved in Section \ref{sec:parity}.
\end{proof}

\begin{theorem}[Uniform hypergeometric Laplace comparison]\label{thm:main_laplace}
Let $N\ge7$ be an integer, $3\le m\le k\le\lfloor N/2\rfloor$, and $N>2m$. If $H\sim\Hyp(N,k,m)$, put $p=k/N$, $M=m(N-m)/(N-1)$ and $a=(N-2m)/(N-2)$. Then for every real $t$,
\begin{equation}\label{eq:main_laplace}
 \log\E e^{t(H-mp)}\le\frac{M}{a^2}\{\log(1-p+pe^{at})-pat\},
\end{equation}
with strict inequality for $t\ne0$. For $k<N/2$,
\begin{equation}\label{eq:main_bern_envelope}
 \sigma^2_{\mathrm{opt}}(H_{N,k,m})
 <\frac{m(N-m)}{N-1}\sigma^2_{\mathrm{opt}}(\Bern(k/N)),
\end{equation}
while for $N$ even and $k=N/2$ equality holds in the optimal-proxy comparison. The same proxy comparison and its equality classification hold for $m=2$ and $2\le k\le\lfloor N/2\rfloor$, $N\ge5$.
\end{theorem}
\begin{proof}
Theorem \ref{odd:thm:envelope} proves the Laplace comparison and the proxy assertions for $m\ge3$; Corollary \ref{cor:two_bern_envelope} treats $m=2$, including the balanced endpoint.
\end{proof}

\begin{theorem}[Universal hypergeometric curvature]\label{thm:curvature}
For all integers $N\ge2$, $1\le k,m<N$, and all $t\in\R$,
\begin{equation}\label{eq:curvature}
 \boxed{\quad v_{N,k,m}(t)\le
 \frac{\min\{k(N-k),\,m(N-m)\}}{4(N-1)}.\quad}
\end{equation}
The bound is sharp over the indicated parameter family. Equality occurs for a singleton marked margin after a suitable exponential tilt, and at zero tilt when both margin products are maximal and $N$ is even.
\end{theorem}

\begin{corollary}[Uniform log-Laplace domination]\label{cor:mgf}
In the same parameter range,
\begin{equation}\label{eq:mgf}
 \log\E e^{t(H-km/N)}\le
 \frac{t^2}{8(N-1)}\min\{k(N-k),m(N-m)\}.
\end{equation}
Consequently the optimal variance proxy obeys
\begin{equation}\label{eq:proxy}
 \sigma^2_{\mathrm{opt}}(H)\le
 \frac{\min\{k(N-k),m(N-m)\}}{4(N-1)}.
\end{equation}
\end{corollary}
\begin{proof}
The centered log-Laplace transform satisfies $\psi(0)=\psi'(0)=0$. Integrate the pointwise second-derivative bound twice, with orientation reversed when $t<0$. For $t\ne0$, the integrated inequality is strict: equality would make the nonnegative analytic curvature gap vanish on the interval between $0$ and $t$, and hence on the whole real line, contradicting $v(t)\to0$ at infinity. The proxy statement follows from the definition.
\end{proof}


Write $L_N=2/[(N-1)\log((N+1)/(N-1))]$ for odd $N$, and let $R_{N,1,m}$ denote the normalized proxy $4N\sigma^2_{\mathrm{opt}}(H_{N,(N-1)/2,m})/[m(N-m)]$.
\begin{theorem}[Uniform fixed-sample odd extremizers]\label{thm:main_uniform}
For each fixed $m\ge2$, all sufficiently large odd $N$ satisfy
\[
 \max_{1\le k<N}\frac{4N}{m(N-m)}\sigma^2_{\mathrm{opt}}(H_{N,k,m})
 =R_{N,1,m}<L_N.
\]
Only the success counts $(N-1)/2$ and $(N+1)/2$ maximize the left side. Their proxy and maximizing tilt admit the fourth- and third-order expansions in Theorem \ref{thm:odd_expansion}.
\end{theorem}
\begin{proof}
The extremizer statement is proved in Theorem \ref{thm:uniform_fixed}; the expansions are given in Theorem \ref{thm:odd_expansion}.
\end{proof}


\begin{corollary}[Weighted finite-exchangeability concentration]\label{cor:tail}
Let $X_1,\ldots,X_N\in[-1,1]$ be exchangeable, $1\le n\le N$, $w\in\R^n$, and let $\widetilde w$ be the $N$-vector obtained by padding $w$ with zeros. Put
\[
 P=I-\frac1N\one\one^\top,\qquad V=\|P\widetilde w\|_2^2.
\]
Then for every $\lambda\in\R$,
\[
 \E\exp\Bigl\{\lambda\sum_{i=1}^nw_i(X_i-\bar X_N)\Bigr\}
 \le\exp\left\{\frac{\lambda^2C_N^\star V}{2}\right\}.
\]
If $V>0$, then for $\delta\in(0,1)$,
\[
 \Pr\left\{\sum_{i=1}^nw_i(X_i-\bar X_N)\ge
 \sqrt{2C_N^\star V\log\frac1\delta}\right\}\le\delta.
\]
Here $C_N^\star$ is given by \eqref{eq:CNexact} and is optimal for every $N$.
\end{corollary}

\section{The two-level reduction and the variational identity}\label{sec:exchangeable}
We include the structural reduction of Lee and Kim \cite[Proposition 2 and Remark 1]{LeeKim}. Let $e_k$ denote the $k$th elementary symmetric polynomial, with $e_j=0$ for indices outside the degree range and $e_0=1$.

\begin{proposition}[Two-level reduction; Lee--Kim]\label{prop:two_level}
Fix $1\le k<N$ and $\rho\ge0$. Every maximizer of
\[
 y\longmapsto e_k(e^{y_1},\ldots,e^{y_N})
 \quad\text{on}\quad
 \left\{y\in\R^N:\sum_i y_i=0,\ \sum_i y_i^2=\rho^2\right\}
\]
has at most two distinct coordinate values.
\end{proposition}
\begin{proof}
Compactness gives a maximizer. The claims for $\rho=0$ and $N=2$ are immediate. Suppose that a maximizer has three distinct coordinates $x_1,x_2,x_3$. Fix all other coordinates $R$ and set $s=x_1+x_2+x_3$, $q=x_1^2+x_2^2+x_3^2$. On the circle defined by these two constraints, the part of the objective depending on the triple is
\[
 A\sum_{i=1}^3e^{x_i}+B\sum_{i=1}^3e^{-x_i},
 \qquad A=e_{k-1}(e^R),\quad B=e^s e_{k-2}(e^R).
\]
Indeed, the terms involving none or all three coordinates are constant, and $e^{x_1+x_2}+e^{x_1+x_3}+e^{x_2+x_3}=e^s\sum_i e^{-x_i}$. The degree range $1\le k<N$ gives $A,B\ge0$ and $A+B>0$.

Put $\phi(x)=Ae^x+Be^{-x}$. At the constrained maximum, Lagrange multipliers give $\phi'(x_i)=\alpha+\beta x_i$. Define
\[
 h(x)=\phi'(x)-\alpha-\beta x,\qquad p(x)=\prod_{i=1}^3(x-x_i),\qquad
 \eta_i=\frac1{p'(x_i)}.
\]
The identities $\sum_i\eta_i=\sum_i x_i\eta_i=0$ show that $\eta$ is tangent to the constraint circle. The second-order necessary condition is
\begin{equation}\label{eq:twolevel_second}
 \sum_{i=1}^3\frac{h'(x_i)}{p'(x_i)^2}\le0.
\end{equation}
Consider the Hermite interpolation polynomial
\[
 H(x)=\sum_{i=1}^3h'(x_i)(x-x_i)
 \left[\frac{p(x)}{(x-x_i)p'(x_i)}\right]^2.
\]
It satisfies $H(x_i)=h(x_i)=0$ and $H'(x_i)=h'(x_i)$. Repeated Rolle's theorem applied to these three double zeros gives a point $\xi$ between the smallest and largest $x_i$ such that
\[
 5!\sum_{i=1}^3\frac{h'(x_i)}{p'(x_i)^2}
 =H^{(5)}(\xi)=h^{(5)}(\xi)=Ae^\xi+Be^{-\xi}>0.
\]
This contradicts \eqref{eq:twolevel_second}.
\end{proof}

\begin{proposition}[Exact variational identity; Lee--Kim]\label{prop:variational}
The constant in \eqref{eq:exchange-intro} satisfies
\begin{equation}\label{eq:LKidentity}
 \boxed{\quad C_N^\star=
 \max_{1\le k,m<N}\frac{8N}{m(N-m)}
 \sup_{t\ne0}\frac{\psi_{N,k,m}(t)}{t^2}.\quad}
\end{equation}
\end{proposition}
\begin{proof}
For $a\perp\one$, average $\exp(\lambda\langle a,x\rangle)$ over coordinate permutations and denote the resulting convex function on $[-1,1]^N$ by $\Phi(x)$. Exchangeability gives $\E e^{\lambda\langle a,X\rangle}=\E\Phi(X)$. Convexity bounds this by the largest value at a cube vertex. A vertex with $k$ positive coordinates, uniformly permuted, is the $k$th Hamming slice. If $S_k$ is a uniform $k$-subset and $y_i=2\lambda a_i$, its log-MGF is
\[
 F_{N,k}(y)=\log\E\exp\Bigl(\sum_{i\in S_k}y_i\Bigr)
 =\log\frac{e_k(e^{y_1},\ldots,e^{y_N})}{\binom Nk}.
\]
For $\lambda a\ne0$, Proposition \ref{prop:two_level} bounds this objective by that of a vector of the same norm with values $(N-m)t/N$ on $m$ coordinates and $-mt/N$ on the others. Its slice sum and squared norm are
\[
 \sum_{i\in S_k}y_i=t(H_{N,k,m}-km/N),\qquad
 \|y\|_2^2=\frac{m(N-m)}N t^2.
\]
Since $\|y\|_2^2=4\lambda^2\|a\|_2^2$, the normalized log-MGF is at most the right side of \eqref{eq:LKidentity}. Constant slices and $\lambda a=0$ contribute zero.

Conversely, for any $k,m,t$, take $X$ uniform on the $k$th Hamming slice and let $a_i=(N-m)t/(2N)$ on a fixed $m$-subset and $a_i=-mt/(2N)$ otherwise. Then
\[
 \langle a,X\rangle=t(H_{N,k,m}-km/N),\qquad
 \|a\|_2^2=\frac{m(N-m)}{4N}t^2.
\]
The defining inequality at $\lambda=1$ gives the reverse bound after taking the supremum over $t\ne0$ and the maximum over $k,m$.
\end{proof}

Applying Corollary \ref{cor:mgf} in \eqref{eq:LKidentity} gives
\begin{equation}\label{eq:globalC}
 C_N^\star\le\frac{N}{N-1}
 \max_{k,m}\frac{\min\{k(N-k),m(N-m)\}}{m(N-m)}
 =\frac{N}{N-1}.
\end{equation}
Theorem \ref{odd:thm:odd} supplies the exact odd value. Together with the sharp examples below, these results prove Theorem \ref{thm:exchange}. Substituting the projected, zero-padded weights into \eqref{eq:exchange-intro} proves Corollary \ref{cor:tail}; its tail bound follows by minimizing the Chernoff exponent.

\section{Parity, maximal variance, and explicit sharp contrasts}\label{sec:parity}
\begin{proposition}[The variance constant and the parity gap]\label{prop:variance_parity}
For $N\ge2$, define the variance constant over nonzero contrasts by
\[
 D_N=\sup_{\substack{X\in[-1,1]^N\ \mathrm{exchangeable}\\a\ne0,\ a\perp\one}}
 \frac{\Var\langle a,X\rangle}{\|a\|_2^2}.
\]
Then
\begin{equation}\label{eq:variance_constant}
 D_N=\frac{4\lfloor N^2/4\rfloor}{N(N-1)}
 =\begin{cases}N/(N-1),&N\text{ even},\\(N+1)/N,&N\text{ odd}.
 \end{cases}
\end{equation}
For even $N$, $C_N^\star=D_N$. For odd $N$, $D_N<C_N^\star<N/(N-1)$, and, with $x=1/N$,
\begin{align}
 C_N^\star-D_N
 &=\frac{2\int_0^x t\operatorname{artanh}t\,dt}
 {(1-x)\operatorname{artanh}x}>0,\label{eq:variance_gap_integral}\\
 &=\frac{2}{3N^2}+\frac{2}{3N^3}+\frac{26}{45N^4}+O(N^{-5}).\label{eq:variance_gap_expansion}
\end{align}
\end{proposition}
\begin{proof}
Exchangeability and $a\perp\one$ give
\[
 \frac{\Var\langle a,X\rangle}{\|a\|_2^2}
 =\E(X_1^2-X_1X_2)
 =\frac1{N-1}\E\sum_{i=1}^N(X_i-\bar X_N)^2.
\]
The sum $\sum_i(x_i-\bar x)^2$ is a convex function on the cube. At a vertex with $k$ positive coordinates it equals $4k(N-k)/N$, whose maximum is $4\lfloor N^2/4\rfloor/N$. A uniformly permuted nearest-balanced vertex attains the maximum for every nonzero $a\perp\one$. This proves \eqref{eq:variance_constant}.

For even $N$, \eqref{eq:globalC} and the variance lower bound prove equality. For odd $N$, use the value established in Theorem \ref{odd:thm:odd}:
\[
 C_N^\star-D_N
 =\frac{x-(1-x^2)\operatorname{artanh}x}
 {(1-x)\operatorname{artanh}x}.
\]
The numerator has derivative $2x\operatorname{artanh}x$ and vanishes at zero, proving \eqref{eq:variance_gap_integral}. Expanding $\operatorname{artanh}x$ gives \eqref{eq:variance_gap_expansion}; $\operatorname{artanh}x>x$ gives $C_N^\star<N/(N-1)$.
\end{proof}

\begin{proposition}[A sharp single-coordinate contrast]\label{prop:explicit_extremizer}
Let $X$ be uniform on the Hamming slice with $k=\lfloor N/2\rfloor$ positive coordinates, and put $a=e_1-N^{-1}\one$. Then
\[
 \|a\|_2^2=\frac{N-1}{N},\qquad
 \langle a,X\rangle=X_1-\bar X_N=2(B-p),
 \quad B\sim\Bern(p),\quad p=\frac{k}{N}.
\]
For odd $N$, the optimal bound \eqref{eq:exchange-intro} is attained at the nonzero argument $\lambda_\star=\log((N+1)/(N-1))$. For even $N$, its sharp coefficient is attained in the limit $\lambda\to0$.
\end{proposition}
\begin{proof}
For odd $N$, write $\alpha=\log((N+1)/(N-1))=\log((1-p)/p)$. The identity $1-p+pe^{2\alpha}=e^\alpha$ gives
\[
 \log\E e^{\alpha\langle a,X\rangle}
 =(1-2p)\alpha=\frac\alpha N
 =\frac{\alpha^2 L_N}{2}\frac{N-1}{N}.
\]
For even $N$, $2(B-p)$ is a symmetric sign, so
\[
 \lim_{\lambda\to0}\frac{2\log\E e^{\lambda\langle a,X\rangle}}
 {\lambda^2\|a\|_2^2}
 =\lim_{\lambda\to0}\frac{2N\log\cosh\lambda}{(N-1)\lambda^2}
 =\frac{N}{N-1}.
\]
\end{proof}

The Bernoulli optimal proxy, used below in both normalizations, is
\begin{equation}\label{eq:bernproxy}
 \sigma^2_{\mathrm{opt}}(\Bern(p))=
 \begin{cases}\displaystyle\frac{1-2p}{2\log((1-p)/p)},&p\ne1/2,\\[5pt]1/4,&p=1/2.
 \end{cases}
\end{equation}
A calculus proof appears in Appendix \ref{app:bernoulli}; see also \cite{ArbelMarchalNguyen}. The sharp examples prove the odd lower bound
\begin{equation}\label{eq:oddL}
 C_N^\star\ge L_N=\frac{2}{(N-1)\log((N+1)/(N-1))},
\end{equation}
while \eqref{eq:globalC} gives the upper bound for all $N$. Since $L_N=1+N^{-1}+O(N^{-2})$, the first-order limit \eqref{eq:asymptotic} follows. The odd upper bound and the complete hypergeometric equality classification are proved in Section \ref{sec:odd_uniform}.


\section{Symmetries, margins, and the polynomial normalization}\label{sec:reductions}
Four complement and transposition identities are fundamental:
\begin{align}
 H_{N,k,m}&\stackrel d=H_{N,m,k},\label{eq:sym1}\\
 H_{N,N-k,m}&\stackrel d=m-H_{N,k,m},\label{eq:sym2}\\
 H_{N,k,N-m}&\stackrel d=k-H_{N,k,m},\label{eq:sym3}\\
 H_{N,N-k,N-m}&\stackrel d=N-k-m+H_{N,k,m}.\label{eq:sym4}
\end{align}
For exponential-tilt variances, reflection $H\mapsto c-H$ simply replaces $t$ by $-t$. Both $k(N-k)$ and $m(N-m)$ remain invariant under their respective complementation, while the minimum of these products is invariant under exchanging $k$ and $m$. Therefore it suffices to prove the curvature theorem under the normalization
\begin{equation}\label{eq:param}
 k=j,\qquad m=j+h,\qquad N=2j+h+r,
 \qquad j\ge1,\ h,r\in\mathbb Z_{\ge0}.
\end{equation}
Here $j=\min\{k,N-k,m,N-m\}$, and the target upper bound becomes
\begin{equation}\label{eq:Bdef}
 B=\frac{j(N-j)}{4(N-1)}.
\end{equation}
Indeed, $j\le m\le N-j$ implies $j(N-j)\le m(N-m)$ by concavity of $x(N-x)$. Both orders of $h$ and $r$ are included, so the normalization covers both signs of the tilt.

\subsection{A generating polynomial with integer coefficients}
For $z>0$, let
\begin{equation}\label{eq:Q}
 Q_j(z)=\sum_{i=0}^j q_i z^i,
 \qquad
 q_i=\binom ji (j+h)_{\underline i}
 \frac{(r+j)!}{(r+i)!},
\end{equation}
where $(a)_{\underline i}=a(a-1)\cdots(a-i+1)$ denotes a falling factorial (with empty product one). We reserve $(a)_i=a(a+1)\cdots(a+i-1)$ for the rising Pochhammer symbol used in hypergeometric functions. Every $q_i$ is strictly positive. The factor independent of $i$ linking $q_i$ to \eqref{eq:hypergeom} cancels in derivatives of $\log Q_j$; hence
\begin{equation}\label{eq:tiltuv}
 u(z)=z\frac{Q_j'(z)}{Q_j(z)}=\E_{\log z}H,
 \qquad
 v(z)=z\frac{d}{dz}u(z)=\Var_{\log z}(H).
\end{equation}

Direct coefficient manipulation gives
\begin{equation}\label{eq:coeffratio}
\frac{q_{i+1}}{q_i}=
\frac{(j-i)(j+h-i)}{(i+1)(r+i+1)}.
\end{equation}
Comparing coefficients of $z^i$ in the following relation now proves the hypergeometric differential equation
\begin{equation}\label{eq:hyperode}
 z(1-z)Q_j''(z)+
 \bigl[r+1+(j+m-1)z\bigr]Q_j'(z)-jm Q_j(z)=0.
\end{equation}
Equation \eqref{eq:hyperode} is a polynomial identity valid also at $z=1$.

\section{Jacobi representation and the essential admissibility bound}\label{sec:jacobi}
The polynomial $Q_j$ is a terminating hypergeometric function,
\[
 Q_j(z)=(r+1)_j\,{}_2F_1(-j,-m;r+1;z).
\]
Pfaff's transformation and the standard Jacobi representation yield the exact identity
\begin{equation}\label{eq:jacobi}
 \boxed{\ Q_j(z)=j!(1-z)^j
 P_j^{(r,h)}\!\left(\frac{1+z}{1-z}\right).\ }
\end{equation}
The factor $(1-z)^j$ cancels the denominator in the Jacobi composition, giving the polynomial \eqref{eq:Q} at every $z$, including $z=1$. Identity \eqref{eq:jacobi} is also verified directly by the coefficient ratio \eqref{eq:coeffratio} and the constant term $Q_j(0)=(r+j)!/r!$ (the rising factorial $(r+1)_j$).

For $r,h\ge0$, the Jacobi polynomial $P_j^{(r,h)}$ has $j$ simple roots in $(-1,1)$, a classical consequence of its orthogonality with respect to the positive Jacobi weight. If these roots are $\xi_1,\ldots,\xi_j$, then the roots of $Q_j$ are $-\alpha_i$, where
\begin{equation}\label{eq:alphas}
 \alpha_i=\frac{1-\xi_i}{1+\xi_i}>0.
\end{equation}
Thus
\begin{equation}\label{eq:factor}
 Q_j(z)=Q_j(0)\prod_{i=1}^j(1+z/\alpha_i).
\end{equation}

\begin{lemma}[Poisson--binomial admissibility]\label{lem:PB}
Let $u=u(z)$ and $v=v(z)$ be defined in \eqref{eq:tiltuv}. For every $z>0$,
\begin{equation}\label{eq:PB}
 \boxed{\quad 0<u<j,\qquad 0<v\le\frac{u(j-u)}j.\quad}
\end{equation}
\end{lemma}
\begin{proof}
Under the exponential tilt $z=e^t$, the normalized probability generating function factors as
\[
 \frac{Q_j(zy)}{Q_j(z)}=
 \prod_{i=1}^j\bigl[1-p_i(z)+p_i(z)y\bigr],
 \qquad p_i(z)=\frac{z}{z+\alpha_i}\in(0,1).
\]
The tilted $H$ is consequently equal in law to a sum of $j$ independent Bernoulli variables with success probabilities $p_1(z),\ldots,p_j(z)$. Therefore $u=\sum p_i$, $v=\sum p_i(1-p_i)=u-\sum p_i^2$. Cauchy--Schwarz gives $\sum p_i^2\ge u^2/j$, establishing \eqref{eq:PB}. Strict positivity of $u$ and $v$ follows from $0<p_i<1$.
\end{proof}

\begin{remark}[Admissible stationary pairs]
Real-rootedness gives the mean--variance constraint \eqref{eq:PB}. At a stationary tilt, the pair $(u,v)$ belongs to the intersection of this region and the algebraic curve derived below.
\end{remark}

\section{The Riccati equation and stationary variance geometry}\label{sec:stationary}
We pass to the tilt parameter $t=\log z$. Write $u(t)=\E_tH$, $v(t)=\Var_t(H)$ and $w(t)=dv(t)/dt$. Then $du/dt=v$. Dividing \eqref{eq:hyperode} by $Q_j(z)$, and using \eqref{eq:tiltuv}, gives the identity
\begin{equation}\label{eq:riccati}
 \boxed{\ (1-z)(v+u^2)+\bigl[r+(j+m)z\bigr]u-jmz=0.\ }
\end{equation}
The equation holds at $z=1$ by continuity. In particular, substituting $z=1$ gives the familiar mean $u(1)=jm/N$.

Differentiating \eqref{eq:riccati} with respect to $t$ gives
\begin{align}\label{eq:differentiated}
 &(1-z)w+\bigl[r+2u+z(j+m-2u-1)\bigr]v\notag\\
 &\hspace{35mm}-z(j-u)(m-u)=0.
\end{align}
At a stationary point of the tilted variance, $w=0$. Set
\begin{equation}\label{eq:Dden}
 \mathcal D(u,v)=(j-u)(m-u)-v(j+m-2u-1).
\end{equation}
Then \eqref{eq:differentiated} becomes
\begin{equation}\label{eq:zfromuv}
 z\mathcal D(u,v)=v(r+2u).
\end{equation}
Since $r\ge0$ and $u,v>0$, the right side is positive; in particular $\mathcal D(u,v)>0$ at any stationary point. Eliminating $z$ between \eqref{eq:riccati} and \eqref{eq:zfromuv} yields
\begin{equation}\label{eq:F}
 \boxed{\begin{aligned}
 \mathcal F(u,v)={}&\bigl[\mathcal D-v(r+2u)\bigr](v+u^2)
 +r\mathcal D\,u\\
 &+(j+m)v(r+2u)u-jmv(r+2u)=0.
 \end{aligned}}
\end{equation}
This eliminates the exponential tilt exactly.

\begin{lemma}[The admissible stationary branch]\label{lem:branch}
For each fixed $0<u<j$, $\mathcal F(u,v)$ is a quadratic polynomial in $v$ with leading coefficient $-(N-1)$ and constant term
\begin{equation}\label{eq:F0}
 \mathcal F(u,0)=u(j-u)(r+u)(m-u)>0.
\end{equation}
Consequently it has exactly one positive real zero, denoted $v_+(u)$. At any stationary point of the actual tilted variance, $v=v_+(u)$. Moreover, for any $B>0$, if $\mathcal F(u,B)\le0$ then $v_+(u)\le B$.
\end{lemma}
\begin{proof}
Collecting powers of $v$ in \eqref{eq:F} gives the asserted leading coefficient and constant term. Since $\mathcal F(u,0)>0$ and its leading coefficient is negative, the two real roots have opposite signs. The positive root is unique. The polynomial is positive strictly between its two roots and negative to the right of the positive root, proving the last assertion.
\end{proof}

Thus a global curvature proof reduces to excluding $v_+(u)>B$ within the admissibility region \eqref{eq:PB}. The polynomial comparison below applies uniformly in $j$.
\section{A universal polynomial positivity principle}\label{sec:positive}
For the stationary comparison, define
\begin{equation}\label{eq:letters}
 T=N-1,\quad L=N-j-1,\quad S=h+r=L-j+1,\quad D=h-r,
\end{equation}
so $T=L+j$, $L\ge j-1$, and $|D|\le S$. The target curvature coefficient is
\[
 B=\frac{j(L+1)}{4T}.
\]
At a stationary mean, introduce
\begin{equation}\label{eq:AandE}
 x=2u-j,\qquad A=j^2(j-1)-Tx^2,
 \qquad E=L(L+1)-2j.
\end{equation}

\begin{lemma}[Universal stationary comparison]\label{lem:universal}
Suppose $j\ge3$, $L\ge j-1$, $S=L-j+1$, and $|D|\le S$. If $|x|<j\sqrt{(j-1)/T}$, where $T=L+j$, then
\begin{equation}\label{eq:masterpositive}
 \boxed{\quad A(D-x)^2+(L+1)(2jDx+Ex^2)\ge0.\quad}
\end{equation}
The inequality holds throughout the stated real parameter domain.
\end{lemma}
\begin{proof}
The condition on $x$ says $A>0$. Furthermore $E\ge0$, because
\[
 E=L(L+1)-2j\ge j(j-1)-2j=j(j-3)\ge0.
\]
If $xD\ge0$, then each of the two terms in \eqref{eq:masterpositive} is nonnegative. We may therefore assume $xD<0$. Write $q=|x|>0$ and $d=|D|\in(0,S]$. The expression to be bounded becomes
\begin{equation}\label{eq:g}
 g(d)=A(d+q)^2+(L+1)(Eq^2-2jqd).
\end{equation}
We split according to the size of $A$.

\smallskip\noindent\emph{Case I: $A\ge j^2/L$.}
Completing the square gives
\begin{align}\label{eq:square1}
 g(d)&=A\left(d+q-\frac{j(L+1)q}{A}\right)^2\notag\\
 &\quad+(L+1)q^2\left[E+2j-\frac{j^2(L+1)}{A}\right].
\end{align}
Here $E+2j=L(L+1)$. By $A\ge j^2/L$, the bracket in the second line is nonnegative. Thus $g(d)\ge0$.

\smallskip\noindent\emph{Case II: $0<A<j^2/L$.}
Since $A=j^2(j-1)-Tq^2$, this assumption implies
\begin{equation}\label{eq:qmin}
 q>q_0:=j\sqrt{\frac{(j-1)L-1}{LT}}.
\end{equation}
We first record the bound
\begin{equation}\label{eq:qbound}
 q_0\ge\frac jL.
\end{equation}
Indeed \eqref{eq:qbound} is equivalent to $(j-1)L^2-2L-j\ge0$. The left side increases with $L\ge j-1\ge2$ and, at $L=j-1$, equals
\[
 (j-1)^3-3j+2=j^2(j-3)+1>0.
\]
Thus $q\ge j/L$ in Case II.

Differentiate the quadratic \eqref{eq:g} with respect to $d$:
\[
 g'(d)=2A(d+q)-2j(L+1)q.
\]
For $0\le d\le S$,
\begin{align*}
 A(d+q)&\le \frac{j^2}{L}(S+q)\\
 &\le j(L+1)q.
\end{align*}
The second inequality follows from $q\ge j/L$ and
\[
 S=L-j+1\le L+1-\frac jL.
\]
Consequently $g$ is nonincreasing on $[0,S]$ and $g(d)\ge g(S)$. Evaluating at $d=S$,
\begin{equation}\label{eq:gS}
 g(S)=A(S+q)^2+(L+1)q(Eq-2jS).
\end{equation}
The final comparison is
\begin{equation}\label{eq:Eq}
 Eq_0\ge2jS.
\end{equation}
Both sides are nonnegative, since $E\ge0$. Squaring \eqref{eq:Eq} and using the definition of $q_0$ shows that it is equivalent to
\begin{equation}\label{eq:Kdef}
 \mathcal K_j(L,S):=
 E^2\bigl((j-1)L-1\bigr)-4S^2L(L+j)\ge0.
\end{equation}
Put
\[
 a=j-3\ge0,\qquad L=j-1+S=a+2+S.
\]
An exact polynomial identity gives
\begin{align}\label{eq:Kpositive}
 \mathcal K_j(L,S)={}&a^2(a+1)(a+3)^3\\
 &+a(a+3)^2(5a^2+16a+10)S\notag\\
 &+(10a^4+80a^3+213a^2+202a+35)S^2\notag\\
 &+(10a^3+64a^2+117a+52)S^3\notag\\
 &+(5a^2+22a+19)S^4+(a+2)S^5.\notag
\end{align}
Every coefficient displayed here is nonnegative for $a,S\ge0$. Therefore \eqref{eq:Kdef} holds for every $j\ge3$ and every $S\ge0$. This proves \eqref{eq:Eq}. Since $q>q_0$, we obtain $Eq\ge2jS$. Equation \eqref{eq:gS} now gives $g(S)\ge0$, and hence $g(d)\ge0$. Both cases establish \eqref{eq:masterpositive}.
\end{proof}

\begin{remark}[An all-degree coefficient identity]
The expression in \eqref{eq:Kpositive} is a polynomial in $a=j-3$ and $S$ with nonnegative coefficients. Its six powers of $S$ establish the comparison on the entire quadrant $a,S\ge0$.
\end{remark}

\subsection{Connecting the positive expression to the stationary equation}
Substitution in \eqref{eq:F} gives the following identity.
\begin{lemma}[Main algebraic identity]\label{lem:identity}
With the notation \eqref{eq:param}, \eqref{eq:letters}, and \eqref{eq:AandE}, and with $B=j(L+1)/(4T)$,
\begin{equation}\label{eq:exactidentity}
 \boxed{\begin{aligned}
 -16T\,\mathcal F\left(\frac{j+x}{2},B\right)
 & =A(D-x)^2\\
 &\quad+(L+1)(2jDx+Ex^2).
 \end{aligned}}
\end{equation}
The identity holds as a rational-polynomial identity in the unrestricted indeterminates $j,h,r,x$ whenever $T\ne0$.
\end{lemma}
\begin{proof}
Insert $m=j+h$ and $N=2j+h+r$ into \eqref{eq:F}, substitute $u=(j+x)/2$ and $v=j(N-j)/(4(N-1))$, and multiply by $16(N-1)$. Group the resulting polynomial by powers of $D=h-r$. Its coefficient of $D^2$ is $-A$ before changing the overall sign; the coefficient of $D$ is $2x(A-j(L+1))$; and the $D$-free part is $-Ax^2-(L+1)Ex^2$. These three coefficients coincide with expansion of the right side after the overall sign. The coefficient form is recorded in Appendix \ref{app:algebra}.
\end{proof}

\section{Proof of the curvature theorem}\label{sec:proof}
\subsection{The unrestricted inner layers \texorpdfstring{$j\ge3$}{j >= 3}}
\begin{proposition}[All-degree curvature domination]\label{prop:allinner}
For every $j\ge3$, every $h,r\ge0$, and $N,m$ given by \eqref{eq:param},
\begin{equation}\label{eq:allinner}
 \sup_{t\in\R}\Var_t(H_{N,j,j+h})\le\frac{j(N-j)}{4(N-1)}.
\end{equation}
\end{proposition}
\begin{proof}
Write $v(t)=\Var_t(H)$ and let $B=j(N-j)/[4(N-1)]$. Since all $q_i$ in \eqref{eq:Q} are strictly positive, the tilted distribution converges to a point mass at $0$ as $t\to-\infty$ and to a point mass at $j$ as $t\to+\infty$. Thus $v(t)\to0$ at either endpoint. The continuous, strictly positive function $v$ attains its global maximum at some finite $t_*$; at that point $w(t_*)=v'(t_*)=0$.

Suppose, for contradiction, that $v_*:=v(t_*)>B$. Let $u_*=u(t_*)$. By the real-root mean--variance inequality \eqref{eq:PB},
\[
 B<v_*\le\frac{u_*(j-u_*)}{j}.
\]
Writing $x_*=2u_*-j$ and recalling $T=N-1$, we obtain
\[
 x_*^2<j^2\frac{j-1}{T},\qquad
 A_*=j^2(j-1)-Tx_*^2>0.
\]
Lemma \ref{lem:universal} applies to $x=x_*$ and $D=h-r$, because $L=N-j-1\ge j-1$ and $|D|\le h+r=L-j+1$. Lemma \ref{lem:identity} consequently gives
\[
 \mathcal F(u_*,B)\le0.
\]
On the other hand $w(t_*)=0$, so the stationary elimination in \eqref{eq:F} shows $\mathcal F(u_*,v_*)=0$. Lemma \ref{lem:branch} states that this zero is the unique positive root of the concave quadratic and that $\mathcal F(u_*,B)\le0$ forces $v_*\le B$, a contradiction. Hence $\sup_t v(t)\le B$.
\end{proof}

\subsection{The singleton layer \texorpdfstring{$j=1$}{j = 1}}
If the smallest margin is one, after the symmetry normalization we may take $k=1$. Then $H$ is a Bernoulli random variable with success probability $m/N$. Under exponential tilt it is still Bernoulli, so $v(t)=p_t(1-p_t)\le1/4$. The right side of \eqref{eq:curvature} is exactly $k(N-k)/[4(N-1)]=1/4$. This proves the case $j=1$ for every $N$.

\subsection{The two-mark layer \texorpdfstring{$j=2$}{j = 2}}
For the two-mark layer, suppose $k=2$ and $2\le m\le N-2$; then $N\ge4$. The three probabilities of $H=0,1,2$, after a common positive factor is removed, have a quadratic generating polynomial $a+2bz+cz^2$ with $a,b,c>0$. Make the change $y=z\sqrt{c/a}$ and define
\[
 \rho=\frac b{\sqrt{ac}}
 =\sqrt{\frac{m(N-m)}{(m-1)(N-m-1)}}.
\]
The polynomial is proportional to $1+2\rho y+y^2$. Since $2\le m,N-m$, the elementary bounds
\[
 \frac{N}{N-2}\le\rho\le2
\]
hold. The upper bound follows by multiplying $m/(m-1)\le2$ and $(N-m)/(N-m-1)\le2$. For the lower bound, observe that
\[
 \rho^2=1+\frac{N-1}{(m-1)(N-m-1)},
 \quad (m-1)(N-m-1)\le\frac{(N-2)^2}4.
\]

Put $y=e^s$. For the shifted three-point law on $\{-1,0,1\}$, elementary differentiation gives
\begin{equation}\label{eq:3pointvar}
 v(s)=\frac{\rho\cosh s+1}{(\rho+\cosh s)^2}.
\end{equation}
For $q=\cosh s\ge1$,
\[
 \frac{d}{dq}\frac{\rho q+1}{(\rho+q)^2}
 =\frac{\rho^2-\rho q-2}{(\rho+q)^3}\le0
\]
because $\rho\le2$. Therefore
\[
 v(s)\le v(0)=\frac1{1+\rho}
 \le\frac{N-2}{2(N-1)}=
 \frac{2(N-2)}{4(N-1)}.
\]
This is the desired symmetric curvature envelope for $j=2$. The cases with another margin equal to one were handled in the preceding subsection.

\subsection{Completion of Theorem \ref{thm:curvature}}
Every admissible $(N,k,m)$ is carried by the complement/transposition symmetries \eqref{eq:sym1}--\eqref{eq:sym4} to the normal form \eqref{eq:param}, where $j$ equals the smallest margin. The case $j=1$ follows from Bernoulli variance; the case $j=2$ from \eqref{eq:3pointvar}; and all $j\ge3$ from Proposition \ref{prop:allinner}. The right side of \eqref{eq:curvature} is invariant under every transformation. This proves Theorem \ref{thm:curvature}.\qed

\section{Parity, maximal variance, and explicit sharp contrasts}\label{sec:parity}
\begin{proposition}[The variance constant and the parity gap]\label{prop:variance_parity}
For $N\ge2$, define the variance constant over nonzero contrasts by
\[
 D_N=\sup_{\substack{X\in[-1,1]^N\ \mathrm{exchangeable}\\a\ne0,\ a\perp\one}}
 \frac{\Var\langle a,X\rangle}{\|a\|_2^2}.
\]
Then
\begin{equation}\label{eq:variance_constant}
 D_N=\frac{4\lfloor N^2/4\rfloor}{N(N-1)}
 =\begin{cases}N/(N-1),&N\text{ even},\\(N+1)/N,&N\text{ odd}.
 \end{cases}
\end{equation}
For even $N$, $C_N^\star=D_N$. For odd $N$, $D_N<C_N^\star<N/(N-1)$, and, with $x=1/N$,
\begin{align}
 C_N^\star-D_N
 &=\frac{2\int_0^x t\operatorname{artanh}t\,dt}
 {(1-x)\operatorname{artanh}x}>0,\label{eq:variance_gap_integral}\\
 &=\frac{2}{3N^2}+\frac{2}{3N^3}+\frac{26}{45N^4}+O(N^{-5}).\label{eq:variance_gap_expansion}
\end{align}
\end{proposition}
\begin{proof}
Exchangeability and $a\perp\one$ give
\[
 \frac{\Var\langle a,X\rangle}{\|a\|_2^2}
 =\E(X_1^2-X_1X_2)
 =\frac1{N-1}\E\sum_{i=1}^N(X_i-\bar X_N)^2.
\]
The sum $\sum_i(x_i-\bar x)^2$ is a convex function on the cube. At a vertex with $k$ positive coordinates it equals $4k(N-k)/N$, whose maximum is $4\lfloor N^2/4\rfloor/N$. A uniformly permuted nearest-balanced vertex attains the maximum for every nonzero $a\perp\one$. This proves \eqref{eq:variance_constant}.

For even $N$, \eqref{eq:globalC} and the variance lower bound prove equality. For odd $N$, use the value established in Theorem \ref{odd:thm:odd}:
\[
 C_N^\star-D_N
 =\frac{x-(1-x^2)\operatorname{artanh}x}
 {(1-x)\operatorname{artanh}x}.
\]
The numerator has derivative $2x\operatorname{artanh}x$ and vanishes at zero, proving \eqref{eq:variance_gap_integral}. Expanding $\operatorname{artanh}x$ gives \eqref{eq:variance_gap_expansion}; $\operatorname{artanh}x>x$ gives $C_N^\star<N/(N-1)$.
\end{proof}

\begin{proposition}[A sharp single-coordinate contrast]\label{prop:explicit_extremizer}
Let $X$ be uniform on the Hamming slice with $k=\lfloor N/2\rfloor$ positive coordinates, and put $a=e_1-N^{-1}\one$. Then
\[
 \|a\|_2^2=\frac{N-1}{N},\qquad
 \langle a,X\rangle=X_1-\bar X_N=2(B-p),
 \quad B\sim\Bern(p),\quad p=\frac{k}{N}.
\]
For odd $N$, the optimal bound \eqref{eq:exchange-intro} is attained at the nonzero argument $\lambda_\star=\log((N+1)/(N-1))$. For even $N$, its sharp coefficient is attained in the limit $\lambda\to0$.
\end{proposition}
\begin{proof}
For odd $N$, write $\alpha=\log((N+1)/(N-1))=\log((1-p)/p)$. The identity $1-p+pe^{2\alpha}=e^\alpha$ gives
\[
 \log\E e^{\alpha\langle a,X\rangle}
 =(1-2p)\alpha=\frac\alpha N
 =\frac{\alpha^2 L_N}{2}\frac{N-1}{N}.
\]
For even $N$, $2(B-p)$ is a symmetric sign, so
\[
 \lim_{\lambda\to0}\frac{2\log\E e^{\lambda\langle a,X\rangle}}
 {\lambda^2\|a\|_2^2}
 =\lim_{\lambda\to0}\frac{2N\log\cosh\lambda}{(N-1)\lambda^2}
 =\frac{N}{N-1}.
\]
\end{proof}

The Bernoulli optimal proxy, used below in both normalizations, is
\begin{equation}\label{eq:bernproxy}
 \sigma^2_{\mathrm{opt}}(\Bern(p))=
 \begin{cases}\displaystyle\frac{1-2p}{2\log((1-p)/p)},&p\ne1/2,\\[5pt]1/4,&p=1/2.
 \end{cases}
\end{equation}
A calculus proof appears in Appendix \ref{app:bernoulli}; see also \cite{ArbelMarchalNguyen}. The sharp examples prove the odd lower bound
\begin{equation}\label{eq:oddL}
 C_N^\star\ge L_N=\frac{2}{(N-1)\log((N+1)/(N-1))},
\end{equation}
while \eqref{eq:globalC} gives the upper bound for all $N$. Since $L_N=1+N^{-1}+O(N^{-2})$, the first-order limit \eqref{eq:asymptotic} follows. The odd upper bound and the complete hypergeometric equality classification are proved in Section \ref{sec:odd_uniform}.


\section{The exact two-draw benchmark}\label{sec:twodraw}
For two draws with $2\le k\le N-2$, the hypergeometric law has three support points. The optimal-proxy formula is the specialization of Atouani, Marchal, and Arbel \cite[Theorem 3.2 and Corollary 3.3]{AtouaniMarchalArbel} under a shift of the support. We give its calculus proof and then optimize the resulting expression over the hypergeometric success count.

\begin{lemma}[Three-mass proxy formula; Atouani--Marchal--Arbel]\label{lem:threemass}
Let $X$ take values $0,1,2$ with strictly positive masses $p_0,p_1,p_2$, and suppose
\[
 \frac{p_1}{2\sqrt{p_0p_2}}\le2.
\]
Then, with the continuous interpretation when $p_0=p_2$,
\begin{equation}\label{eq:threeproxy}
 \sigma^2_{\mathrm{opt}}(X)=\frac{2(p_0-p_2)}{\log(p_0/p_2)}.
\end{equation}
\end{lemma}
\begin{proof}
By reflection take $p_0\ge p_2$ and set $A=p_1$, $B=2\sqrt{p_0p_2}$, $a=\frac12\log(p_0/p_2)$. Writing $f(x)=\log(A+B\cosh x)$ gives
\[
 \E e^{tX}=e^t\{A+B\cosh(t-a)\},\qquad
 f'(a)=p_0-p_2.
\]
For $x>0$,
\[
 f'''(x)=\frac{B\sinh x\,[A^2-AB\cosh x-2B^2]}{(A+B\cosh x)^3}\le0,
\]
since $A\le2B$. Thus $f'(x)/x$ decreases on $(0,\infty)$, and $f$ is even. For $a>0$, integration on either side of $a$ yields
\[
 f(x)-f(a)\le\frac{f'(a)}{2a}(x^2-a^2),\qquad x\in\mathbb R.
\]
After substituting $x=t-a$ and centering $X$, this gives the optimal quadratic bound with equality at $t=2a$. When $a=0$, the same derivative argument gives the optimal coefficient $f''(0)=2p_0$, equal to the continuous limit of \eqref{eq:threeproxy}.
\end{proof}

\begin{proposition}[All two-draw laws for an odd population]\label{prop:odd2}
Let $N\ge5$ be odd. Among all $H\sim\Hyp(N,k,2)$, the largest normalized optimal variance proxy equals
\begin{equation}\label{eq:Ltwo}
 L_N^{(2)}=\frac{4}{(N-2)\log\frac{N+1}{N-3}},
\end{equation}
attained at $k=(N-1)/2$ or $k=(N+1)/2$. For every such $N$,
\begin{equation}\label{eq:twostrict}
 L_N^{(2)}<L_N,
 \qquad L_N-L_N^{(2)}=\frac{2}{3N^3}+O(N^{-4}).
\end{equation}
\end{proposition}
\begin{proof}
For $2\le k\le N-2$, the three masses of $H$ are
\[
 p_0=\frac{(N-k)(N-k-1)}{N(N-1)},\quad
 p_1=\frac{2k(N-k)}{N(N-1)},\quad
 p_2=\frac{k(k-1)}{N(N-1)}.
\]
The hypothesis of Lemma \ref{lem:threemass} follows from
$\bigl(p_1/(2\sqrt{p_0p_2})\bigr)^2=k(N-k)/[(k-1)(N-k-1)]\in[1,4]$.
Consequently, for $k\ne N/2$,
\[
 \sigma^2_{\mathrm{opt}}(H)
 =\frac{2(N-2k)}{N\log\frac{(N-k)(N-k-1)}{k(k-1)}}.
\]
On replacing $k$ by $N-k$ if necessary, put $x=N-2k>0$. Then
\[
 \log\frac{(N-k)(N-k-1)}{k(k-1)}
 =2\operatorname{artanh}(x/N)+2\operatorname{artanh}(x/(N-2)).
\]
As $\operatorname{artanh}(x/b)/x$ is strictly increasing for $0<x<b$, the normalized proxy is largest when $x=1$. The endpoint laws $k=1,N-1$ are Bernoulli; Hoeffding's lemma bounds their normalized contributions by $N/[2(N-2)]<1$, whereas the central two-draw contribution exceeds one. This proves \eqref{eq:Ltwo}.

Set $y=N-1$. A direct integral identity gives
\begin{align*}
 &(y-1)\log\frac{y+2}{y-2}-2y\log\frac{y+2}{y}\\
 &\qquad=2y\int_0^1(1-u)
 \left[\frac1{y^2-(2-u)^2}-\frac1{y^2-u^2}\right]du>0.
\end{align*}
This is equivalent to $L_N^{(2)}<L_N$. The asymptotic difference follows by expanding $\operatorname{artanh}x=x+x^3/3+x^5/5+O(x^7)$ in both exact formulas.
\end{proof}

\begin{corollary}[Two-draw Bernoulli envelope]\label{cor:two_bern_envelope}
For every integer $N\ge5$ and $2\le k\le\lfloor N/2\rfloor$,
\[
 \sigma^2_{\mathrm{opt}}(H_{N,k,2})
 \le\frac{2(N-2)}{N-1}\sigma^2_{\mathrm{opt}}(\Bern(k/N)).
\]
The inequality is strict for $k<N/2$ and is an equality for even $N$ and $k=N/2$.
\end{corollary}
\begin{proof}
For $k=N/2$ the three-mass law is symmetric, and Lemma \ref{lem:threemass} gives its proxy as its variance $2(N-2)/[4(N-1)]$, establishing equality. For $k<N/2$, put $d=1-2k/N>0$ and
\[
 A_1=2\operatorname{artanh}d,\qquad
 B_1=2\operatorname{artanh}\frac{Nd}{N-2}.
\]
The formula in the proof of Proposition \ref{prop:odd2} gives
$\sigma^2_{\mathrm{opt}}(H)=2d/(A_1+B_1)$. Since $\operatorname{artanh}(x)/x$ is strictly increasing on $(0,1)$ and $k\ge2$,
\[
 B_1>\frac{N}{N-2}A_1.
\]
Consequently
\[
 \frac{2d}{A_1+B_1}<\frac{(N-2)d}{(N-1)A_1}
 =\frac{2(N-2)}{N-1}\sigma^2_{\mathrm{opt}}(\Bern(k/N)),
\]
as asserted.
\end{proof}


\section{Uniform Laplace envelopes and odd-population extremizers}\label{sec:odd_uniform}
Write $g(d)=d/\operatorname{artanh}d$ with $g(0)=1$, and recall $L_N=[(N-1)\operatorname{artanh}(1/N)]^{-1}$. The Laplace comparison below applies to both parities; only the final optimization over discrete success counts uses oddness.

\begin{theorem}[Parity-uniform Bernoulli envelope]\label{odd:thm:envelope}
Let $N\ge7$ be an integer with $3\le m\le k\le\lfloor N/2\rfloor$ and $N>2m$, and let $H\sim\Hyp(N,k,m)$. Write $\psi=\psi_{N,k,m}$ and set
\[
 p=\frac{k}{N},\quad q=1-p,\quad d=1-2p,\quad
 M=\frac{m(N-m)}{N-1},\quad a=\frac{N-2m}{N-2}>0.
\]
For every $t\in\R$,
\begin{equation}\label{odd:eq:laplace}
 \boxed{\quad
 \psi(t)\le\frac{M}{a^2}\bigl[\log(q+pe^{at})-pat\bigr].\quad}
\end{equation}
The Laplace inequality is strict for $t\ne0$. If $k<N/2$, then
\begin{equation}\label{odd:eq:BernEnvelope}
 \sigma^2_{\rm opt}(H)<M\sigma^2_{\rm opt}(\Bern(p)),\qquad
 \frac{4N}{m(N-m)}\sigma^2_{\rm opt}(H)
 <\frac{N}{N-1}g(d).
\end{equation}
If $N$ is even and $k=N/2$, the two optimal proxies are equal: both sides of the first comparison equal $M/4$, and both sides of the normalized comparison equal $N/(N-1)$. For odd $N$, $d\ge 1/N$, so the last expression in \eqref{odd:eq:BernEnvelope} is at most $L_N$.
All these estimates are finite-population statements.
\end{theorem}

\subsection{Tilted mean and variance as a curve}
Throughout the envelope proof assume the hypotheses of Theorem~\ref{odd:thm:envelope}.
Let $\E_t$ and $\Var_t$ denote exponential tilting by $e^{tH}$, and write
\[
 u(t)=\E_tH,\qquad v(t)=\Var_t(H).
\]
Every mass on $\{0,\ldots,m\}$ is positive. Hence $u'(t)=v(t)>0$ and $u$ maps $\R$ bijectively onto $(0,m)$. Thus $v$ defines a smooth function $V(u)$. At the endpoints $V(u)\to0$, and, writing $\mu=mp$,
\begin{equation}\label{odd:eq:variance0}
 V(\mu)=Mpq=\frac M4(1-d^2).
\end{equation}

\begin{lemma}[Real-root admissibility]\label{odd:lem:PB}
For every $0<u<m$,
\begin{equation}\label{odd:eq:PB}
 0<V(u)\le\frac{u(m-u)}m.
\end{equation}
\end{lemma}
\begin{proof}
Apply the Jacobi factorization and Bernoulli representation from Section \ref{sec:jacobi} with degree $m$ and parameters $N-k-m$ and $k-m$. If the tilted Bernoulli probabilities are $p_1,\ldots,p_m$, then
\[
 u=\sum_i p_i,\qquad V=\sum_i p_i(1-p_i)
 =u-\sum_i p_i^2\le u-\frac{u^2}{m}.
\]
\end{proof}

Define
\[
 A(u)=u(u+N-k-m),\qquad C(u)=(k-u)(m-u)
\]
and the polynomial
\begin{equation}\label{odd:eq:F}
 F(u,v)=-(N-1)v^2+[A+C-A'C+C'A]v+AC.
\end{equation}
Primes on $A,C$ denote differentiation with respect to $u$.

\begin{lemma}[Differential identity]\label{odd:lem:ode}
Writing $y=u-\mu$, the true variance curve obeys
\begin{equation}\label{odd:eq:curveODE}
 F(u,V(u))+NyV(u)V'(u)=0.
\end{equation}
\end{lemma}
\begin{proof}
Let $z=e^t$. The tilted adjacent-mass relation, summed over the support, gives
\[
 \E_t[H(H+N-k-m)]=z\E_t[(k-H)(m-H)].
\]
Consequently
\[
 z=\frac{A(u)+V(u)}{C(u)+V(u)}.
\]
Both numerator and denominator are positive. Differentiate its logarithm with respect to $u$, using $dt/du=1/V$:
\[
 \frac{A'+V'}{A+V}-\frac{C'+V'}{C+V}=\frac1V.
\]
Multiplication and collection yield $(C-A)VV'=F(u,V)$. Since $C-A=km-Nu=-Ny$, this is \eqref{odd:eq:curveODE}. The identity also holds at $y=0$.
\end{proof}

\subsection{A quadratic envelope and an all-parameter certificate}
Define a quadratic polynomial in $u$ by
\begin{equation}\label{odd:eq:P}
 P(u)=\frac M4(1-d^2)+ad(u-\mu)-\frac{a^2}{M}(u-\mu)^2.
\end{equation}
The envelope depends quadratically on the tilted mean.

\begin{lemma}[Quadratic variance envelope]\label{odd:lem:barrier}
For every $0<u<m$,
\begin{equation}\label{odd:eq:barrier}
 \boxed{\quad V(u)\le P(u),\quad}
\end{equation}
with equality if and only if $u=\mu$.
\end{lemma}

For the variance comparison, put
\[
 z_0=\frac{2(u-\mu)}M,\quad
 G=\frac{N^2(m-1)(N-m-1)}{4(N-2)^2(N-1)}>0
\]
and, for free variables $d,z_0$, define
\begin{equation}\label{odd:eq:W}
 W(d,z_0)=N-4+3Nd^2-4Nad z_0+[(N+1)a^2-M]z_0^2.
\end{equation}
An exact substitution into \eqref{odd:eq:F} gives
\begin{equation}\label{odd:eq:residual}
 \boxed{\quad F(u,P(u))+N(u-\mu)P(u)P'(u)
 =-G(u-\mu)^2W(d,z_0).\quad}
\end{equation}
At a noncentral contact $V=P$, the sign of the residual is determined by $W$.

Set
\begin{align*}
 \delta&=\frac{N-2m}{N},&
 b&=\frac{N-m}{N-2},&
 c&=\frac{(N-m)[3N^2-4N-4m(N-1)]}{(N-1)(N-2)^2},\\
 E(d,z_0)&=1-d^2+4bdz_0-cz_0^2.
\end{align*}
The exact admissibility identity is
\begin{equation}\label{odd:eq:E}
 \frac{u(m-u)}m-P(u)=\frac{m(m-1)}{4(N-1)}E(d,z_0).
\end{equation}
Thus, at any contact $V=P$, Lemma~\ref{odd:lem:PB} implies $E\ge0$.
Also $m\le k\le N/2$ implies $0\le d\le\delta$; equality $d=0$ occurs precisely at the even balanced margin $k=N/2$.

Define
\begin{equation}\label{odd:eq:lambda}
 \lambda=\frac{m^2-mN+3m+N^2-4N}{2(N-m)},\qquad
 \alpha=\frac{N-4-\lambda}{\delta^2},
\end{equation}
and
\[
 A_0=3N+\lambda+\alpha,\qquad
 B_0=-2Na-2\lambda b,\qquad
 C_0=(N+1)a^2-M+\lambda c.
\]
The following is a polynomial identity after clearing the positive denominators:
\begin{equation}\label{odd:eq:certificate}
 \boxed{\quad
 W=\lambda E+\alpha(\delta^2-d^2)
   +A_0d^2+2B_0dz_0+C_0z_0^2.\quad}
\end{equation}
Moreover,
\begin{equation}\label{odd:eq:det}
 A_0C_0-B_0^2
 =\frac{N^2U_0V_0}{(N-2m)^2(N-2)^2(N-1)},
\end{equation}
where
\begin{align*}
 U_0&=N^2-2N-m^2-mN+3m,\\
 V_0&=m^3-6m^2-2mN^2+14mN-3m+N^3-7N^2+6N.
\end{align*}

For a direct all-degree sign verification write
\[
 j=m-3\ge0,\qquad r=N-2m\ge1.
\]
Then
\begin{align}
 2(N-m)\lambda
 &=3j^2+3jr+13j+r^2+5r+12>0,\label{odd:eq:lambda_pos}\\
 2(N-m)(N-4-\lambda)
 &=j^2+3jr+3j+r^2+5r>0,\label{odd:eq:alpha_pos}\\
 U_0&=j^2+3jr+5j+r^2+7r+6>0,\label{odd:eq:U_pos}\\
 V_0&=j^3+4j^2r+3j^2+4jr^2+10jr+r^3+5r^2>0.\label{odd:eq:V_pos}
\end{align}
It follows that $\lambda,\alpha>0$, $A_0>0$, and $A_0C_0-B_0^2>0$.
The quadratic form in \eqref{odd:eq:certificate} is positive definite. Therefore
\begin{equation}\label{odd:eq:Wpositive}
 E\ge0,\quad 0\le d\le\delta\quad\Longrightarrow\quad W(d,z_0)>0.
\end{equation}
The coefficient identities establish strict positivity for every allowed $N,m$.

\begin{proof}[Proof of Lemma~\ref{odd:lem:barrier}]
At the two endpoints the envelope is positive. Explicitly,
\begin{align}
 P(0)&=\frac{km(m-1)[(N-m)(N-2)-k(m-1)]}
 {(N-m)(N-2)^2(N-1)}>0,\label{odd:eq:end0}\\
 P(m)&=\frac{m(N-k)(m-1)[k(m-1)+(N-2m)(N-1)]}
 {(N-m)(N-2)^2(N-1)}>0.\label{odd:eq:endm}
\end{align}
For the first sign, $k\le N-m$ bounds the bracket below by
$(N-m)(N-m-1)>0$. The second sign is immediate.
Let $f(u)=V(u)-P(u)$, extending continuously to $[0,m]$.
Then $f(0)<0$, $f(m)<0$, and \eqref{odd:eq:variance0} gives $f(\mu)=0$.

At any zero $u\ne\mu$ of $f$, we have $P=V>0$. Combining
\eqref{odd:eq:curveODE} with \eqref{odd:eq:residual} gives
\begin{equation}\label{odd:eq:crossing}
 f'(u)=\frac{G(u-\mu)W(d,z_0)}{NP(u)}.
\end{equation}
At this zero, \eqref{odd:eq:E} and \eqref{odd:eq:PB} imply $E\ge0$, so
\eqref{odd:eq:Wpositive} applies. Thus every zero left of $\mu$ has strictly negative derivative, and every zero right of $\mu$ has strictly positive derivative.

If $f$ were positive somewhere in $(0,\mu)$, the left endpoint of a component of $\{f>0\}$ would lie in $(0,\mu)$, and would be a zero with nonnegative derivative, a contradiction. If $f$ were positive somewhere in $(\mu,m)$, the right endpoint of a positive component would give a zero with nonpositive derivative, again a contradiction. Hence $f\le0$ everywhere. Any additional interior zero would be a local maximum of $f$ and have derivative zero, contradicting \eqref{odd:eq:crossing}. This proves \eqref{odd:eq:barrier} and its equality statement.
\end{proof}

\subsection{Integration gives the scaled Bernoulli Laplace bound}
Write $y(t)=u(t)-\mu$. Lemma~\ref{odd:lem:barrier} says
\[
 y'(t)\le \mathcal P(y(t)),\qquad
 \mathcal P(y)=Mpq+ady-\frac{a^2}{M}y^2,
\]
with strict inequality for $t\ne0$. The exact solution starting at zero to
$Y'=\mathcal P(Y)$ is
\begin{equation}\label{odd:eq:Y}
 Y(t)=\frac{M}{a}\left[\frac{pe^{at}}{q+pe^{at}}-p\right].
\end{equation}
For the two-sided comparison, put
\[
 \Delta(t)=\mathcal P(y(t))-y'(t),\qquad
 h(t)=ad-\frac{a^2}{M}(y(t)+Y(t)).
\]
Then $\Delta(t)>0$ for $t\ne0$, and $e(t)=y(t)-Y(t)$ solves
\[
 e'=he-\Delta,\quad e(0)=0,\qquad
 e(t)=-e^{\int_0^t h}\int_0^t e^{-\int_0^s h}\Delta(s)\,ds.
\]
Thus $y(t)<Y(t)$ for $t>0$ and $y(t)>Y(t)$ for $t<0$.
Since $\psi'(t)=y(t)$ and $\psi(0)=0$, integration, with the correct orientation for negative $t$, yields
\[
 \psi(t)<\int_0^tY(s)\,ds
 =\frac{M}{a^2}[\log(q+pe^{at})-pat]\quad(t\ne0).
\]
This proves \eqref{odd:eq:laplace}.

The exact Bernoulli proxy is
\begin{equation}\label{odd:eq:bern}
 \sigma^2_{\rm opt}(\Bern(p))=\frac{q-p}{2\log(q/p)}
 =\frac{d}{4\artanh d}.
\end{equation}
On substituting $s=at$ in the supremum, \eqref{odd:eq:laplace} implies
$\sigma^2_{\rm opt}(H)\le M\sigma^2_{\rm opt}(\Bern(p))$.
To optimize the comparison, extend the hypergeometric quotient continuously to zero. It tends to zero at either infinity, so its maximum is attained at a finite point. At every nonzero tilt it is strictly below the Bernoulli upper envelope. If $d>0$, its value at zero is $Mpq$, strictly below $M\sigma^2_{\rm opt}(\Bern(p))=Md/(4\artanh d)$, because
\[
 \artanh d=\int_0^d\frac{dx}{1-x^2}<\frac{d}{1-d^2}\qquad(0<d<1).
\]
Thus the proxy comparison is strict for $k<N/2$. If $d=0$, the integrated Laplace bound and $\log\cosh x\le x^2/2$ give $\sigma^2_{\rm opt}(H)\le M/4$. The variance at zero tilt is $Mpq=M/4$, so equality holds in the optimal-proxy comparison. Finally, $g(d)$ strictly decreases for $d>0$ by the same integral inequality. In the odd case $d\ge1/N$, and hence $g(d)\le g(1/N)$. This completes Theorem~\ref{odd:thm:envelope} for both parities.

\subsection{Exact odd-population extremizers}
\begin{theorem}[Odd-population hypergeometric extremizers]\label{odd:thm:odd}
For every odd $N\ge3$,
\begin{equation}\label{odd:eq:exact}
 \boxed{\quad
 \max_{1\le k,m<N}\frac{4N}{m(N-m)}\sigma^2_{\mathrm{opt}}(H_{N,k,m})
 =\frac{1}{(N-1)\operatorname{artanh}(1/N)}.\quad}
\end{equation}
Equality holds precisely for $m\in\{1,N-1\}$ and $k\in\{(N-1)/2,(N+1)/2\}$.
\end{theorem}
\begin{proof}
Complementation of either margin replaces the centered variable by its negative. It leaves the optimal proxy unchanged, and sample complementation preserves $m(N-m)$. Reduce first to $1\le k,m\le(N-1)/2$.

For $m=1$, the Bernoulli formula \eqref{eq:bernproxy} gives
\[
 \frac{4N}{N-1}\sigma^2_{\mathrm{opt}}(H_{N,k,1})
 =\frac{N}{N-1}g(1-2k/N)\le L_N,
\]
with equality exactly at $k=(N-1)/2$. For $2\le m\le k$, Corollary \ref{cor:two_bern_envelope} and Theorem \ref{odd:thm:envelope} give strict inequality.

If $k<m$, transposition gives
\[
 \frac{4N}{m(N-m)}\sigma^2_{\mathrm{opt}}(H_{N,k,m})
 =\frac{k(N-k)}{m(N-m)}
 \frac{4N}{k(N-k)}\sigma^2_{\mathrm{opt}}(H_{N,m,k})<L_N,
\]
because the first factor is strictly below one and the transposed law is in the canonical triangle. Restoring the complements gives the stated equality cases. For $N=3$ only the singleton layer occurs; for $N=5$ the singleton and two-draw layers exhaust the canonical triangle.
\end{proof}

Applying \eqref{eq:LKidentity} proves the odd formula in Theorem \ref{thm:exchange}. The comparison is finite-population and covers every sample size, including sequences with $m_N\to\infty$, whether $m_N=o(N)$ or $m_N/N$ stays bounded away from zero.


\section{Optimal proxies for nearly balanced odd populations}\label{sec:odd_asym}
The strict singleton comparison of Theorem \ref{odd:thm:odd} has a small gap for nearly balanced populations. This section determines that gap through fourth order and the location of the globally maximizing tilt through third order, for every fixed sample size and fixed success imbalance. The cases $m=1$ and $m=2$ agree with the Bernoulli formula \eqref{eq:bernproxy} and the three-point specialization in Section \ref{sec:twodraw}. The theorem gives the expansion for arbitrary fixed $m$, together with uniqueness of the global optimizer and analytic remainders.

Let $N$ be odd, let $s\ge1$ be a fixed odd integer, and put
\begin{equation}\label{eq:oddparams}
 k=\frac{N-s}{2},\qquad H_N\sim\Hyp(N,k,m),\qquad
 R_{N,s,m}=\frac{4N}{m(N-m)}\sigma^2_{\mathrm{opt}}(H_N).
\end{equation}
Here $m\ge1$ is fixed while $N\to\infty$ runs through odd integers larger than $s+2m$. The prescribed centering gives a Bernoulli-success fraction $p=(1-s/N)/2$. The regime includes the most nearly balanced success count when $s=1$.

\begin{theorem}[Fourth-order sharp-proxy expansion]\label{thm:odd_expansion}
Fix $m\ge1$ and a positive odd integer $s$. For all sufficiently large odd $N$, the function
\[
 t\longmapsto\frac{2}{t^2}\log\E\exp\{t(H_N-mk/N)\},
\]
extended continuously at $t=0$, has a unique global maximizer $t_{N,s,m}>0$. The maximizer and the optimal normalized proxy have the expansions
\begin{align}
 t_{N,s,m}&=\frac{4s}{N}+\frac{4s(m-1)}{N^2}
 +\frac{4s(6m+s^2-6)}{3N^3}+O_{m,s}(N^{-4}),\label{eq:tiltexp}\\
 R_{N,s,m}
 &=1+\frac1N+\frac{1-s^2/3}{N^2}
 +\frac{1-(2m-1)s^2/3}{N^3}\notag\\
 &\quad+\frac{1-(4m^2-2m-1)s^2/3-4s^4/45}{N^4}
 +O_{m,s}(N^{-5}).\label{eq:Rexp}
\end{align}
The remainder is analytic in $N^{-1}$ in a neighborhood of zero for every fixed $m,s$.
\end{theorem}

\begin{corollary}[Fourth-order singleton dominance]\label{cor:fixedm_gap}
For every fixed integer $m\ge2$, let $H_N\sim\Hyp(N,(N-1)/2,m)$, $N$ odd, and let
\[
 L_N=\frac{1}{(N-1)\operatorname{artanh}(1/N)}.
\]
Then
\begin{equation}\label{eq:fixedgap}
 \boxed{\quad L_N-R_{N,1,m}
 =\frac{2(m-1)}{3N^3}
 +\frac{2(m-1)(2m+1)}{3N^4}+O_m(N^{-5}).\quad}
\end{equation}
In particular, $R_{N,1,m}<L_N$ for all sufficiently large odd $N$, for every fixed $m\ge2$. More generally, if $s\ge3$ is fixed and odd, then
\begin{equation}\label{eq:offs_gap}
 L_N-R_{N,s,m}=\frac{s^2-1}{3N^2}+O_{m,s}(N^{-3})>0
\end{equation}
for all sufficiently large odd $N$.
\end{corollary}

\subsection{Cumulants from factorial moments}\label{sec:cumulants_odd}
We compute the cumulants entering the globally optimized expansion in Theorem \ref{thm:odd_expansion}. Write $\varepsilon=N^{-1}$. For $r\le m$, the factorial moments of $H_N$ satisfy
\begin{equation}\label{eq:factorial_moments}
 \E(H_N)_{\underline r}
 =\frac{(m)_{\underline r}(k)_{\underline r}}
 {(N)_{\underline r}},\qquad k=\frac{N-s}{2}.
\end{equation}
The ordinary moments up to order six are obtained by the Stirling transformation
\[
 \E H_N^q=\sum_{r=0}^q\left\{\begin{matrix}q\\r\end{matrix}\right\}
 \E(H_N)_{\underline r}.
\]
The expressions remain valid when $r>m$, with zero falling factorials. Let $\kappa_j(\varepsilon)$ denote the $j$th cumulant of $H_N$. Elementary simplification of the first six moments yields
\begin{align}
 \kappa_2(e)&=\frac m4\left[1-(m-1)e-(m+s^2-1)e^2
 +(m-1)(s^2-1)(e^3+e^4)\right]+O_{m,s}(e^5),\label{eq:kappa2exp}\\
 \kappa_3(e)&=\frac{ms}{4}\left[e-3(m-1)e^2
 +(2m^2-9m-s^2+7)e^3\right]+O_{m,s}(e^4),\label{eq:kappa3exp}\\
 \kappa_4(e)&=-\frac m8+\frac{m(m-1)}2e
 +\frac{m(-6m^2+19m+4s^2-13)}8e^2+O_{m,s}(e^3),\label{eq:kappa4exp}\\
 \kappa_5(e)&=-\frac{ms}{2}e+O_{m,s}(e^2),\qquad
 \kappa_6(e)=\frac m4+O_{m,s}(e).\label{eq:kappa56exp}
\end{align}
Here $e=\varepsilon$. All expansions follow from factorial moments through order six and the cumulant recursion
\[
 \kappa_r=\mu_r-\sum_{j=1}^{r-1}\binom{r-1}{j-1}\kappa_j\mu_{r-j},
 \qquad \mu_r=\E H^r.
\]
For the variance and third cumulant one can verify the first two identities directly from
\begin{align}
 \kappa_2(\varepsilon)
 &=\frac m4(1-s^2\varepsilon^2)\frac{1-m\varepsilon}{1-\varepsilon},\label{eq:exactk2}\\
 \kappa_3(\varepsilon)
 &=\kappa_2(\varepsilon)\,s\varepsilon
 \frac{1-2m\varepsilon}{1-2\varepsilon}.
 \label{eq:exactk3}
\end{align}
Formula \eqref{eq:kappa4exp} follows from \eqref{eq:factorial_moments} for $r=1,2,3,4$, using
\[
 \kappa_4=\E(H-\E H)^4-3\kappa_2^2.
\]
All these identities are rational in $\varepsilon$ for fixed $m,s$ and are valid in a neighborhood of the origin. The zero-parameter limit is $H_0\sim\mathrm{Binomial}(m,1/2)$, for which odd cumulants vanish and $\kappa_4(0)=-m/8$.

\subsection{Localization and uniqueness of the maximizing tilt}
Let
\[
 Q_\varepsilon(t)=\begin{cases}
  2\psi_\varepsilon(t)/t^2,&t\ne0,\\
  \kappa_2(\varepsilon),&t=0.
 \end{cases}
\]
For every fixed $m,s$, the finite probability generating polynomial depends analytically on $(t,\varepsilon)$ near $(0,0)$, and its value at $(0,0)$ is one. Consequently $Q_\varepsilon$ is jointly real analytic in a neighborhood of $(0,0)$. At $\varepsilon=0$,
\begin{equation}\label{eq:q0}
 Q_0(t)=\frac{2m}{t^2}\log\cosh(t/2),\qquad Q_0(0)=m/4.
\end{equation}
The strict inequality $\log\cosh x<x^2/2$ for $x\ne0$ implies
\begin{equation}\label{eq:qstrict}
 Q_0(t)<Q_0(0)\qquad(t\ne0).
\end{equation}
Furthermore, $|H_N-\E H_N|\le m$ implies $0\le Q_\varepsilon(t)\le2m/|t|$ for $t\ne0$, uniformly in the admissible populations. On every compact interval, $Q_\varepsilon$ converges uniformly to $Q_0$, including at zero by analytic continuation. Combining this compact convergence with the uniform tail bound shows that every global maximizer of $Q_\varepsilon$ tends to zero as $\varepsilon\downarrow0$.

Near $(0,0)$, the cumulant expansion gives
\begin{equation}\label{eq:Qexpand}
 Q_\varepsilon(t)=\kappa_2(\varepsilon)
 +\frac{\kappa_3(\varepsilon)}3t
 +\frac{\kappa_4(\varepsilon)}{12}t^2
 +\frac{\kappa_5(\varepsilon)}{60}t^3
 +\frac{\kappa_6(\varepsilon)}{360}t^4+O_{m,s}(t^5).
\end{equation}
At the limiting law, $\kappa_5(0)=0$, and
\[
 \partial_t^2Q_0(0)=\frac{\kappa_4(0)}6=-\frac m{48}<0.
\]
Thus $Q_\varepsilon$ is strictly concave on some fixed neighborhood of zero for all sufficiently small $\varepsilon$. The derivative at zero is $\kappa_3(\varepsilon)/3>0$ for $\varepsilon>0$ small. By the analytic implicit-function theorem, there exists a unique critical point $t(\varepsilon)$ in that neighborhood; it is analytic in $\varepsilon$, and the localization argument establishes that it is the unique global maximizer. Differentiation of the critical point equation at zero gives
\[
 t'(0)=-\frac{\partial_\varepsilon\partial_t Q(0,0)}{\partial_t^2Q(0,0)}
 =-\frac{ms/12}{-m/48}=4s.
\]
This proves the leading term of \eqref{eq:tiltexp} and the positivity of the maximizing tilt. The next subsection computes the remaining coefficients.

\subsection{Evaluation through fourth order}
The analytic critical-point equation can be expanded coefficient by coefficient:
\[
 0=\frac{\kappa_3(e)}3+\frac{\kappa_4(e)}6t
 +\frac{\kappa_5(e)}{20}t^2+\frac{\kappa_6(e)}{90}t^3+O_{m,s}(t^4).
\]
Set $t=a e+b e^2+c e^3+O(e^4)$. The coefficients of $e,e^2,e^3$, in order, give
\[
 a=4s,\qquad b=4s(m-1),\qquad c=\frac{4s(6m+s^2-6)}3.
\]
The terms omitted from the derivative contribute $O_{m,s}(e^4)$ because the cumulants are analytic and $t=O(e)$. This establishes the full tilt expansion in \eqref{eq:tiltexp}.

For the optimal value through order $e^4$, it suffices to use
\[
 \widehat t=4se+4s(m-1)e^2.
\]
Indeed $t_* -\widehat t=O(e^3)$ and $Q_e'(t_*)=0$, so Taylor's theorem around the true optimizer gives $Q_e(t_*)-Q_e(\widehat t)=O(e^6)$. With $\kappa_5(0)=0$ and $\kappa_6(0)=m/4$, substitution into \eqref{eq:Qexpand} yields
\begin{align*}
 \sigma^2_{\mathrm{opt}}(H_N)
 &=\kappa_2+\frac{\kappa_3}3\widehat t
 +\frac{\kappa_4}{12}\widehat t^2
 +\frac{\kappa_5}{60}\widehat t^3
 +\frac{\kappa_6}{360}\widehat t^4+O_{m,s}(e^5)\\
 &=\frac m4\Bigl[1-(m-1)e+(1-m-s^2/3)e^2\\
 &\hspace{14mm}-(m-1)(1+s^2/3)e^3\\
 &\hspace{14mm}+\{1-m-(2m^2-m-1)s^2/3-4s^4/45\}e^4\Bigr]
 +O_{m,s}(e^5).
\end{align*}
Multiplying by $4/[m(1-me)]$ proves the fourth-order proxy expansion \eqref{eq:Rexp}. The analytic global optimizer supplies the Taylor remainder for each fixed $m,s$.

The exact singleton benchmark is
\begin{equation}\label{eq:singletonL}
 L_N=\frac{1}{(N-1)\operatorname{artanh}(1/N)}
 =1+e+\frac23e^2+\frac23e^3+\frac{26}{45}e^4+O(e^5).
\end{equation}
Subtracting \eqref{eq:Rexp} from this expression proves Corollary \ref{cor:fixedm_gap}.

\subsection{Finite-size interpretation}
The $N^{-1}$ coefficient in \eqref{eq:Rexp} is one. The $N^{-2}$ coefficient records the success imbalance through $-s^2/3$ and is common to all fixed sample sizes. For $s=1$, the first sample-size contribution is the order-$N^{-3}$ term $-2(m-1)/3$ relative to the singleton law. Thus \eqref{eq:fixedgap} quantifies the separation between the optimized proxies. The maximizing tilt has leading term $4s/N$.

\section{Uniform extremizers at fixed sample size}\label{sec:uniform_fixed}
At each fixed sample size, the following results maximize the proxy over every success count and compute its local separation from the Bernoulli envelope.

\subsection{A two-parameter analytic family}
Fix an integer $m\ge1$. Write $e=N^{-1}$, $d=1-2k/N$, $p=(1-d)/2$, and $q=1-p$. For $i=0,\ldots,m$, define
\begin{equation}\label{eq:analytic_masses}
 P_i(e,d)=\binom mi
 \frac{\prod_{r=0}^{i-1}(p-re)\prod_{r=0}^{m-i-1}(q-re)}
 {\prod_{r=0}^{m-1}(1-re)}.
\end{equation}
Empty products equal one. For physical parameters this is the hypergeometric probability mass function. Vandermonde's falling-factorial identity gives $\sum_iP_i=1$ and $\sum_i iP_i=mp$. In a rectangle around $(e,d)=(0,0)$, every factor is positive and the masses are analytic. At $e=0$ the law is $\mathrm{Binomial}(m,p)$. Define
\[
 \psi(e,d,t)=\log\sum_{i=0}^mP_i(e,d)e^{t(i-mp)},\qquad
 Q(e,d,t)=\frac{2\psi(e,d,t)}{t^2},
\]
with $Q(e,d,0)=\partial_t^2\psi(e,d,0)$. Division by $t^2$ is removable because $\psi(e,d,0)=\partial_t\psi(e,d,0)=0$. Let
\[
 \mathcal R_m(e,d)=\frac{4}{m(1-me)}\max_{t\in\mathbb R}Q(e,d,t).
\]

\begin{lemma}[Joint analytic optimizer]\label{lem:joint_optimizer}
For each fixed $m$, there is a rectangle around $(0,0)$ on which $Q(e,d,\cdot)$ has a unique global maximizer $\tau_m(e,d)$. Both $\tau_m$ and $\mathcal R_m$ are real analytic. Moreover,
\begin{align}
 \tau_m(e,-d)&=-\tau_m(e,d),&
 \mathcal R_m(e,-d)&=\mathcal R_m(e,d),\label{eq:joint_symmetry}\\
 \tau_m(e,0)&=0,&
 \mathcal R_m(e,0)&=\frac1{1-e}.\label{eq:balanced_exact}
\end{align}
\end{lemma}
\begin{proof}
At $(e,d)=(0,0)$, $Q(t)=2m\log\cosh(t/2)/t^2$ has its unique maximum at zero and $Q_{tt}(0)=-m/48$. For nearby parameters, $|H-mp|\le m$ gives $0\le Q(t)\le2m/|t|$. Choose $T$ so large that this tail bound lies strictly below $m/4$. On $[-T,T]$, joint continuity, including at $t=0$, yields uniform convergence to the limiting objective. Its strict gap outside any neighborhood of zero forces every maximizing point into that neighborhood. Joint analyticity and $Q_{tt}(0,0,0)<0$ then give a fixed smaller neighborhood where $Q_{tt}<0$. The analytic implicit-function theorem applied to $Q_t=0$ gives one critical point there, which the localization argument identifies as the unique global maximum.

The identity $P_i(e,-d)=P_{m-i}(e,d)$ implies $Q(e,-d,t)=Q(e,d,-t)$, proving \eqref{eq:joint_symmetry}. At $d=0$ the objective is even in $t$, so uniqueness gives $\tau_m(e,0)=0$. Its value is the exact variance
\[
 \partial_t^2\psi(e,0,0)=\frac m4\frac{1-me}{1-e},
\]
which proves \eqref{eq:balanced_exact}.
\end{proof}

\begin{theorem}[Uniform fixed-sample maximum]\label{thm:uniform_fixed}
For each fixed integer $m\ge1$, there exists $N_0(m)$ such that, for every odd $N\ge N_0(m)$,
\begin{equation}\label{eq:fixed_all_k}
 \max_{1\le k<N}\frac{4N}{m(N-m)}
 \sigma^2_{\mathrm{opt}}(H_{N,k,m})=R_{N,1,m}.
\end{equation}
The maximizing success counts are exactly $k=(N-1)/2$ and $k=(N+1)/2$. If $m\ge2$, this maximum is strictly below $L_N$. Consequently, for every fixed $M\ge2$, all sample sizes $2\le m\le M$ and all success counts are simultaneously below $L_N$ once $N$ is sufficiently large.
\end{theorem}
\begin{proof}
At $e=0$, the Bernoulli proxy formula gives
\[
 \mathcal R_m(0,d)=g(d):=\frac d{\operatorname{artanh}d},\qquad g(0)=1.
\]
In particular $g(d)=1-d^2/3+O(d^4)$. The even analytic function from Lemma \ref{lem:joint_optimizer} therefore has the form
\[
 \mathcal R_m(e,d)=\frac1{1-e}-d^2A_m(e,d^2),\qquad A_m(0,0)=\frac13,
\]
where $A_m$ is analytic. Choose $d_0>0$ and $e_0>0$ inside its analytic rectangle such that
\[
 A_m(e,d^2)+d^2\partial_2A_m(e,d^2)>0
 \quad(0\le e\le e_0,\ 0\le d\le d_0).
\]
It follows that $\partial_d\mathcal R_m(e,d)<0$ for $0<d\le d_0$. Thus, among odd-population success counts in this interval, the maximum is uniquely attained at $d=1/N=e$.

For the remaining success counts, the elementary-symmetric-mean inequality gives
\begin{equation}\label{eq:maclaurin_mgf}
 \E e^{tH_{N,k,m}}\le(q+pe^t)^m.
\end{equation}
Indeed, $\E e^{tH}$ is the normalized $m$th elementary symmetric polynomial of $N$ positive weights, $k$ equal to $e^t$ and the others equal to one. At fixed sum, averaging a pair of weights increases that polynomial by $(a-b)^2e_{m-2}(\mathrm{rest})/4\ge0$. Repeatedly averaging the largest and smallest weights decreases the sum of their squares by $(a-b)^2/2$; consequently their difference tends to zero, and the weights converge to their common mean. The elementary symmetric polynomial never decreases during this procedure. The case $m=1$ is equality. This proves \eqref{eq:maclaurin_mgf}. After centering, the Bernoulli proxy formula yields
\[
 \frac{4N}{m(N-m)}\sigma^2_{\mathrm{opt}}(H_{N,k,m})
 \le\frac{g(d)}{1-me}\le\frac{g(d_0)}{1-me}\quad(d\ge d_0).
\]
Since $g(d_0)<1$, this last bound is below one for sufficiently small $e>0$, whereas $\mathcal R_m(e,e)=1+e+O_m(e^2)>1$. Reflection covers negative $d$. This proves \eqref{eq:fixed_all_k} and the equality cases. The strict comparison with $L_N$ follows from Corollary \ref{cor:fixedm_gap}. Taking the largest of finitely many thresholds proves the assertion for $2\le m\le M$.
\end{proof}


\subsection{Exact local separation at finite population size}
The polynomial proof of Theorem \ref{thm:curvature} also applies to the normal form with integer $j\ge1$ and real $h,r\ge0$. In this extension, the coefficients in \eqref{eq:Q} use Gamma functions for factorial ratios. Their ratio, the differential equation, and the stationary identity remain unchanged; Jacobi orthogonality holds for real $h,r\ge0$, and the positivity certificate uses only $S=h+r\ge0$. The cases $j\ge3$ and the two boundary layers therefore prove the same curvature bound. This observation allows a continuous success-count perturbation at $k=N/2$, including odd $N$.

\begin{theorem}[Exact imbalance coefficient]\label{thm:exact_imbalance}
Fix integers $N\ge5$ and $1\le m\le\lfloor N/2\rfloor$. Extend the hypergeometric masses to real $k=N(1-d)/2$ by \eqref{eq:analytic_masses}, with $e=1/N$. For all sufficiently small $|d|$, the optimal normalized proxy has a unique maximizing tilt, analytic in $d$, and
\begin{align}
 \tau_m(1/N,d)&=T_{N,m}d+O_{N,m}(d^3),\label{eq:finite_tilt}\\
 \mathcal R_m(1/N,d)&=\frac N{N-1}-A_{N,m}d^2+O_{N,m}(d^4),\label{eq:finite_imbalance}
\end{align}
where, writing $D_{N,m}=N^2-N-3m(N-m)>0$,
\begin{align}
 T_{N,m}&=\frac{4(N-2m)(N-3)(N-1)}{(N-2)D_{N,m}},\label{eq:finite_T}\\
 A_{N,m}&=\frac N{3(N-1)}+
 \frac{2N^2(N-4)(m-1)(N-m-1)}
 {3(N-2)^2(N-1)D_{N,m}}.\label{eq:finite_A}
\end{align}
For $m\ge2$, $A_{N,m}>A_{N,1}$. Thus the Bernoulli envelope strictly dominates the $m$-draw proxy for every sufficiently small nonzero real imbalance.
\end{theorem}
\begin{proof}
At $d=0$, use the real-margin extension with $j=m$ and $h=r=N/2-m\ge0$. The curvature envelope equals the variance
\[
 B=\frac{m(N-m)}{4(N-1)}.
\]
Its twice-integrated bound is strict at every nonzero $t$, by the analytic-gap argument in Corollary \ref{cor:mgf}; hence $Q(1/N,0,t)$ has its unique global maximum at zero. The masses remain positive for $|d|$ sufficiently small. Their finite support gives uniform tail control, so the localization argument in Lemma \ref{lem:joint_optimizer} applies at this fixed $N$.

Factorial moments give, at $d=0$,
\begin{align*}
 \kappa_2(d)&=B(1-d^2),\\
 \kappa_3(d)&=B\frac{N-2m}{N-2}d+O(d^3),\\
 \kappa_4(0)&=-\frac{BD_{N,m}}{2(N-3)(N-1)}.
\end{align*}
Since $m(N-m)\le N^2/4$, $D_{N,m}\ge N(N-4)/4>0$. Consequently $Q_{tt}(1/N,0,0)=\kappa_4(0)/6<0$, and the analytic implicit-function theorem gives a unique analytic maximizing tilt. If $b=B(N-2m)/(N-2)$ and $c=\kappa_4(0)$, the stationary equation gives $T_{N,m}=-2b/c$, proving \eqref{eq:finite_T}. Substitution into
\[
 Q(d,t)=B(1-d^2)+\frac{bd}{3}t+\frac c{12}t^2
 +O_{N,m}(d^4)\qquad(t=O(d))
\]
gives
\[
 \max_tQ(d,t)=B-\left(B+\frac{b^2}{3c}\right)d^2+O_{N,m}(d^4).
\]
Reflection makes the optimized value even in $d$, so the stated remainder follows. Multiplying by $4N/[m(N-m)]$ yields
\[
 A_{N,m}=\frac N{N-1}\left[
 1-\frac{2(N-2m)^2(N-3)(N-1)}{3(N-2)^2D_{N,m}}\right].
\]
Subtracting $N/[3(N-1)]$ and putting over a common denominator gives exactly the positive second term of \eqref{eq:finite_A}. The singleton proxy is $\mathcal R_1(1/N,d)=N g(d)/(N-1)$, so its quadratic coefficient is $A_{N,1}=N/[3(N-1)]$. The asserted local domination follows from the strict coefficient difference.
\end{proof}

\subsection{Contraction on compact success-proportion intervals}
The first finite-population correction gives a comparison away from the balanced point as well.

\begin{theorem}[Uniform Bernoulli envelope at fixed sample size]\label{thm:compact_contraction}
Fix $m\ge2$ and $0<\eta<1/2$. There exists $N_1(m,\eta)$ such that, for all $N\ge N_1(m,\eta)$ and $\eta\le k/N\le1-\eta$,
\begin{equation}\label{eq:compact_envelope}
 \sigma^2_{\mathrm{opt}}(H_{N,k,m})
 \le\frac{m(N-m)}{N-1}\,
 \sigma^2_{\mathrm{opt}}(\Bern(k/N)).
\end{equation}
Equality holds when $k=N/2$; the inequality is strict otherwise.
\end{theorem}
\begin{proof}
By reflection let $d\in[0,1-2\eta]$. At $e=0$ the binomial proxy has its unique maximizing tilt $t_0(d)=4\operatorname{artanh}d$, with continuous value $t_0(0)=0$. Its second derivative there is negative: for $d>0$, the Bernoulli proof in Appendix \ref{app:bernoulli} gives $G''(t_0)=c-pq>0$ and $Q''(t_0)=-2mG''(t_0)/t_0^2<0$; at $d=0$ the derivative is $-m/48$. Compactness, the uniform tail bound, and the strict uniqueness at $e=0$ allow the localization and implicit-function arguments to be applied uniformly on this interval. Thus $\mathcal R_m(e,d)$ is analytic on a neighborhood of its compact parameter interval for sufficiently small $|e|$.

For $p,q$ in this compact interval, expanding \eqref{eq:analytic_masses} gives
\[
 P_i(e,d)=\binom mi p^iq^{m-i}\left[
 1+\frac e2\left\{m(m-1)-\frac{i(i-1)}p
 -\frac{(m-i)(m-i-1)}q\right\}\right]+O_{m,\eta}(e^2).
\]
If $Z=q+pe^t$, binomial factorial moments therefore yield
\begin{equation}\label{eq:first_order_psi}
 \psi(e,d,t)=m\log Z-mpt+
 \frac{em(m-1)}2\left[1-\frac{pe^{2t}+q}{Z^2}\right]+O_{m,\eta}(e^2)
\end{equation}
on compact $t$ intervals. Analytic divisibility by $t^2$ permits the same expansion for $Q$ at zero. Evaluation at the unique limiting maximizer gives, with $a=\log(q/p)=2\operatorname{artanh}d$,
\[
 \mathcal R_m(e,d)=g(d)+e\left[m g(d)-(m-1)\frac{d^2}{a^2pq}\right]+O_{m,\eta}(e^2).
\]
The continuous interpretation is used at $d=0$. Since $\mathcal R_1(e,d)=g(d)/(1-e)$, their difference vanishes identically at $e=0$ and also at $d=0$ for every small $e$. It is even in $d$. Joint analyticity consequently factors it exactly as
\[
 \mathcal R_1(e,d)-\mathcal R_m(e,d)=e d^2 B_m(e,d),
\]
where $B_m$ is continuous on the compact interval and
\[
 B_m(0,d)=\frac{(m-1)g(d)}{d^2}
 \left[\frac{d}{(1-d^2)\operatorname{artanh}d}-1\right],
 \qquad B_m(0,0)=\frac{2(m-1)}3.
\]
For $d>0$, integration gives $\operatorname{artanh}d<d/(1-d^2)$, so $B_m(0,d)>0$. Its compact minimum is positive. Uniform continuity then gives $B_m(e,d)>0$ for all sufficiently small $e>0$. This proves the normalized comparison, including strictness, and multiplying through gives \eqref{eq:compact_envelope}.
\end{proof}

\begin{remark}[Relation to the finite-population envelope]
Theorem \ref{thm:compact_contraction} gives a local analytic comparison for both parities on compact success-proportion intervals. For odd populations, Theorem \ref{thm:main_laplace} gives the strict proxy comparison at every finite sample size in the canonical triangular domain. Theorem \ref{thm:uniform_fixed} further identifies the maximizing success count at each fixed sample size for all sufficiently large odd $N$.
\end{remark}


\section{Pointwise curvature and its weighted integral}\label{sec:curvature_obstruction}
The exact odd-population constant defines the factor
\[
 \rho_N=\frac{N-1}{N}L_N=\frac{1}{N\operatorname{artanh}(1/N)}<1.
\]
The corresponding pointwise inequality takes the form
\begin{equation}\label{eq:falseoddcurve}
 \psi_{N,k,m}''(t)\le\rho_N\frac{m(N-m)}{4(N-1)}.
\end{equation}
The next proposition gives an exact tilted variance above the right side of \eqref{eq:falseoddcurve}.

\begin{proposition}[Exact exceedance of the constant pointwise bound]\label{prop:obstruction}
For $N=15$, $k=7$, $m=3$, and $t=\log(7/6)$,
\begin{equation}\label{eq:obs_exact}
 \frac{4(N-1)}{m(N-m)}\psi_{N,k,m}''(t)
 =\frac{307887776}{308178025}
 >\frac{675}{676}>\rho_{15}.
\end{equation}
This nearest-balanced population violates \eqref{eq:falseoddcurve} at the stated tilt.
\end{proposition}

\begin{proof}
With $z=7/6$, the tilted hypergeometric weights for $j=0,1,2,3$ are proportional to
\[
 w_j=\binom7j\binom8{3-j}z^j.
\]
They are rational, so both the mean and the variance are rational. Direct summation gives
\[
 \E_tH=\frac{26313}{17555},\qquad
 \operatorname{Var}_t(H)=\frac{197927856}{308178025}.
\]
Since $m(N-m)/[4(N-1)]=9/14$, the first equality in \eqref{eq:obs_exact} follows by multiplying the variance by $14/9$. The first strict inequality follows from
\[
 \frac{307887776}{308178025}-\frac{675}{676}
 =\frac{111969701}{208328344900}>0.
\]
Finally, the power series for $\operatorname{artanh}$ at $x=1/15$ yields
\[
 \operatorname{artanh}(1/15)>\frac1{15}+\frac1{3\cdot15^3},
 \quad\text{hence}\quad
 \rho_{15}<\left(1+\frac1{3\cdot15^2}\right)^{-1}=\frac{675}{676}.
\]
Combining these two inequalities proves \eqref{eq:obs_exact}.
\end{proof}

\begin{proposition}[Integrated curvature identity]\label{prop:integral}
For every centered random variable with finite moment generating function, writing $\psi(t)=\log\E e^{tX}$ gives
\begin{equation}\label{eq:curv_integral}
 \frac{2\psi(t)}{t^2}=2\int_0^1(1-u)\psi''(ut)\,du,
 \qquad t\ne0.
\end{equation}
The proxy objective is therefore a weighted average of curvature along the segment from zero to the chosen tilt.
\end{proposition}
\begin{proof}
Twice integrate $\psi''$ using $\psi(0)=\psi'(0)=0$ and substitute $s=ut$; the formula is valid for either sign of $t$.
\end{proof}

Theorem \ref{odd:thm:envelope} bounds the weighted integral \eqref{eq:curv_integral} through the mean-dependent envelope of Lemma \ref{odd:lem:barrier}. Proposition \ref{prop:obstruction} identifies a point of this variance curve above the constant threshold in \eqref{eq:falseoddcurve}.


\section{Jacobi root dispersion and variance geometry}\label{sec:structure}
\subsection{A root-spread inequality for Jacobi polynomials}
Write $Q_j(z)=\text{const.}\prod_i(z+\alpha_i)$ as in \eqref{eq:factor}. Because
\[
 v(\log z)=\sum_{i=1}^j\frac{z\alpha_i}{(z+\alpha_i)^2}
 =\frac14\sum_{i=1}^j\left[1-\left(\frac{z-\alpha_i}{z+\alpha_i}\right)^2\right],
\]
Theorem \ref{thm:curvature} yields the following root inequality, valid for all Jacobi degrees and all nonnegative integer parameters $h,r$:
\begin{equation}\label{eq:rootineq}
 \boxed{\quad
 \sum_{i=1}^j\left(\frac{z-\alpha_i}{z+\alpha_i}\right)^2
 \ge\frac{j(j-1)}{N-1}
 \quad(z>0),\qquad N=2j+h+r.\quad}
\end{equation}
At $z=1$, \eqref{eq:rootineq} recovers the ordinary hypergeometric variance bound. For all $z>0$, it gives the same lower bound for the sum of squared transformed Jacobi zeros.

\subsection{An exact finite-population variance envelope}
The central hypergeometric variance is
\[
 \Var(H)=\frac{km(N-k)(N-m)}{N^2(N-1)}.
\]
The ratio of this variance to the all-tilt symmetric envelope is
\[
 \frac{4\max\{k(N-k),m(N-m)\}}{N^2}\le1.
\]
Thus the envelope is saturated at zero tilt whenever the larger margin product equals $N^2/4$. For $m=1$ and $k\ne N/2$, the tilted variance is maximized at $t=\log((N-k)/k)\ne0$. Theorem \ref{thm:curvature} bounds the maximizing variance for every pair of margins.

\subsection{A probabilistic interpretation of the positivity argument}
The Poisson--binomial bound $v\le u(j-u)/j$ defines a region in the mean--variance plane. At each stationary tilt, the pair $(u,v)$ also satisfies $\mathcal F(u,v)=0$. For $j\ge3$, Lemmas \ref{lem:universal} and \ref{lem:identity} place the admissible stationary branch at or below $v=B$. The singleton and two-mark calculations complete the curvature bound for all degrees.

\appendix
\section{Factorial-moment and cumulant expansions}\label{app:cumulants}
\subsection{Cumulants from factorial moments}\label{sec:cumulants_odd}
We compute the cumulants entering the globally optimized expansion in Theorem \ref{thm:odd_expansion}. Write $\varepsilon=N^{-1}$. For $r\le m$, the factorial moments of $H_N$ satisfy
\begin{equation}\label{eq:factorial_moments}
 \E(H_N)_{\underline r}
 =\frac{(m)_{\underline r}(k)_{\underline r}}
 {(N)_{\underline r}},\qquad k=\frac{N-s}{2}.
\end{equation}
The ordinary moments up to order six are obtained by the Stirling transformation
\[
 \E H_N^q=\sum_{r=0}^q\left\{\begin{matrix}q\\r\end{matrix}\right\}
 \E(H_N)_{\underline r}.
\]
The expressions remain valid when $r>m$, with zero falling factorials. Let $\kappa_j(\varepsilon)$ denote the $j$th cumulant of $H_N$. Elementary simplification of the first six moments yields
\begin{align}
 \kappa_2(e)&=\frac m4\left[1-(m-1)e-(m+s^2-1)e^2
 +(m-1)(s^2-1)(e^3+e^4)\right]+O_{m,s}(e^5),\label{eq:kappa2exp}\\
 \kappa_3(e)&=\frac{ms}{4}\left[e-3(m-1)e^2
 +(2m^2-9m-s^2+7)e^3\right]+O_{m,s}(e^4),\label{eq:kappa3exp}\\
 \kappa_4(e)&=-\frac m8+\frac{m(m-1)}2e
 +\frac{m(-6m^2+19m+4s^2-13)}8e^2+O_{m,s}(e^3),\label{eq:kappa4exp}\\
 \kappa_5(e)&=-\frac{ms}{2}e+O_{m,s}(e^2),\qquad
 \kappa_6(e)=\frac m4+O_{m,s}(e).\label{eq:kappa56exp}
\end{align}
Here $e=\varepsilon$. All expansions follow from factorial moments through order six and the cumulant recursion
\[
 \kappa_r=\mu_r-\sum_{j=1}^{r-1}\binom{r-1}{j-1}\kappa_j\mu_{r-j},
 \qquad \mu_r=\E H^r.
\]
For the variance and third cumulant one can verify the first two identities directly from
\begin{align}
 \kappa_2(\varepsilon)
 &=\frac m4(1-s^2\varepsilon^2)\frac{1-m\varepsilon}{1-\varepsilon},\label{eq:exactk2}\\
 \kappa_3(\varepsilon)
 &=\kappa_2(\varepsilon)\,s\varepsilon
 \frac{1-2m\varepsilon}{1-2\varepsilon}.
 \label{eq:exactk3}
\end{align}
Formula \eqref{eq:kappa4exp} follows from \eqref{eq:factorial_moments} for $r=1,2,3,4$, using
\[
 \kappa_4=\E(H-\E H)^4-3\kappa_2^2.
\]
All these identities are rational in $\varepsilon$ for fixed $m,s$ and are valid in a neighborhood of the origin. The zero-parameter limit is $H_0\sim\mathrm{Binomial}(m,1/2)$, for which odd cumulants vanish and $\kappa_4(0)=-m/8$.

\section{Analytic contraction on compact success-proportion intervals}\label{app:contraction}
\subsection{Fixed-sample comparison}
The first finite-population correction gives a comparison away from the balanced point as well.

\begin{theorem}[Uniform Bernoulli envelope at fixed sample size]\label{thm:compact_contraction}
Fix $m\ge2$ and $0<\eta<1/2$. There exists $N_1(m,\eta)$ such that, for all $N\ge N_1(m,\eta)$ and $\eta\le k/N\le1-\eta$,
\begin{equation}\label{eq:compact_envelope}
 \sigma^2_{\mathrm{opt}}(H_{N,k,m})
 \le\frac{m(N-m)}{N-1}\,
 \sigma^2_{\mathrm{opt}}(\Bern(k/N)).
\end{equation}
Equality holds when $k=N/2$; the inequality is strict otherwise.
\end{theorem}
\begin{proof}
By reflection let $d\in[0,1-2\eta]$. At $e=0$ the binomial proxy has its unique maximizing tilt $t_0(d)=4\operatorname{artanh}d$, with continuous value $t_0(0)=0$. Its second derivative there is negative: for $d>0$, the Bernoulli proof in Appendix \ref{app:bernoulli} gives $G''(t_0)=c-pq>0$ and $Q''(t_0)=-2mG''(t_0)/t_0^2<0$; at $d=0$ the derivative is $-m/48$. Compactness, the uniform tail bound, and the strict uniqueness at $e=0$ allow the localization and implicit-function arguments to be applied uniformly on this interval. Thus $\mathcal R_m(e,d)$ is analytic on a neighborhood of its compact parameter interval for sufficiently small $|e|$.

For $p,q$ in this compact interval, expanding \eqref{eq:analytic_masses} gives
\[
 P_i(e,d)=\binom mi p^iq^{m-i}\left[
 1+\frac e2\left\{m(m-1)-\frac{i(i-1)}p
 -\frac{(m-i)(m-i-1)}q\right\}\right]+O_{m,\eta}(e^2).
\]
If $Z=q+pe^t$, binomial factorial moments therefore yield
\begin{equation}\label{eq:first_order_psi}
 \psi(e,d,t)=m\log Z-mpt+
 \frac{em(m-1)}2\left[1-\frac{pe^{2t}+q}{Z^2}\right]+O_{m,\eta}(e^2)
\end{equation}
on compact $t$ intervals. Analytic divisibility by $t^2$ permits the same expansion for $Q$ at zero. Evaluation at the unique limiting maximizer gives, with $a=\log(q/p)=2\operatorname{artanh}d$,
\[
 \mathcal R_m(e,d)=g(d)+e\left[m g(d)-(m-1)\frac{d^2}{a^2pq}\right]+O_{m,\eta}(e^2).
\]
The continuous interpretation is used at $d=0$. Since $\mathcal R_1(e,d)=g(d)/(1-e)$, their difference vanishes identically at $e=0$ and also at $d=0$ for every small $e$. It is even in $d$. Joint analyticity consequently factors it exactly as
\[
 \mathcal R_1(e,d)-\mathcal R_m(e,d)=e d^2 B_m(e,d),
\]
where $B_m$ is continuous on the compact interval and
\[
 B_m(0,d)=\frac{(m-1)g(d)}{d^2}
 \left[\frac{d}{(1-d^2)\operatorname{artanh}d}-1\right],
 \qquad B_m(0,0)=\frac{2(m-1)}3.
\]
For $d>0$, integration gives $\operatorname{artanh}d<d/(1-d^2)$, so $B_m(0,d)>0$. Its compact minimum is positive. Uniform continuity then gives $B_m(e,d)>0$ for all sufficiently small $e>0$. This proves the normalized comparison, including strictness, and multiplying through gives \eqref{eq:compact_envelope}.
\end{proof}

\begin{remark}[Relation to the finite-population envelope]
Theorem \ref{thm:main_laplace} now gives the finite-population proxy comparison for both parities in the canonical sub-half-population range, with equality at the balanced even margin. The present analytic argument records the first finite-population correction uniformly over compact success-proportion intervals, and Theorem \ref{thm:uniform_fixed} identifies the maximizing success count for each fixed sample size in sufficiently large odd populations.
\end{remark}

\section{Pointwise curvature and its weighted integral}\label{sec:curvature_obstruction}
The exact odd-population constant defines the factor
\[
 \rho_N=\frac{N-1}{N}L_N=\frac{1}{N\operatorname{artanh}(1/N)}<1.
\]
The corresponding pointwise inequality takes the form
\begin{equation}\label{eq:falseoddcurve}
 \psi_{N,k,m}''(t)\le\rho_N\frac{m(N-m)}{4(N-1)}.
\end{equation}
The next proposition gives an exact tilted variance above the right side of \eqref{eq:falseoddcurve}.

\begin{proposition}[Exact exceedance of the constant pointwise bound]\label{prop:obstruction}
For $N=15$, $k=7$, $m=3$, and $t=\log(7/6)$,
\begin{equation}\label{eq:obs_exact}
 \frac{4(N-1)}{m(N-m)}\psi_{N,k,m}''(t)
 =\frac{307887776}{308178025}
 >\frac{675}{676}>\rho_{15}.
\end{equation}
This nearest-balanced population violates \eqref{eq:falseoddcurve} at the stated tilt.
\end{proposition}

\begin{proof}
With $z=7/6$, the tilted hypergeometric weights for $j=0,1,2,3$ are proportional to
\[
 w_j=\binom7j\binom8{3-j}z^j.
\]
They are rational, so both the mean and the variance are rational. Direct summation gives
\[
 \E_tH=\frac{26313}{17555},\qquad
 \operatorname{Var}_t(H)=\frac{197927856}{308178025}.
\]
Since $m(N-m)/[4(N-1)]=9/14$, the first equality in \eqref{eq:obs_exact} follows by multiplying the variance by $14/9$. The first strict inequality follows from
\[
 \frac{307887776}{308178025}-\frac{675}{676}
 =\frac{111969701}{208328344900}>0.
\]
Finally, the power series for $\operatorname{artanh}$ at $x=1/15$ yields
\[
 \operatorname{artanh}(1/15)>\frac1{15}+\frac1{3\cdot15^3},
 \quad\text{hence}\quad
 \rho_{15}<\left(1+\frac1{3\cdot15^2}\right)^{-1}=\frac{675}{676}.
\]
Combining these two inequalities proves \eqref{eq:obs_exact}.
\end{proof}

\begin{proposition}[Integrated curvature identity]\label{prop:integral}
For every centered random variable with finite moment generating function, writing $\psi(t)=\log\E e^{tX}$ gives
\begin{equation}\label{eq:curv_integral}
 \frac{2\psi(t)}{t^2}=2\int_0^1(1-u)\psi''(ut)\,du,
 \qquad t\ne0.
\end{equation}
The proxy objective is therefore a weighted average of curvature along the segment from zero to the chosen tilt.
\end{proposition}
\begin{proof}
Twice integrate $\psi''$ using $\psi(0)=\psi'(0)=0$ and substitute $s=ut$; the formula is valid for either sign of $t$.
\end{proof}

Theorem \ref{odd:thm:envelope} bounds the weighted integral \eqref{eq:curv_integral} through the mean-dependent envelope of Lemma \ref{odd:lem:barrier}. Proposition \ref{prop:obstruction} identifies a point of this variance curve above the constant threshold in \eqref{eq:falseoddcurve}.

\section{Coefficient forms of the stationary identities}\label{app:algebra}
The stationary polynomial \eqref{eq:F} expands as
\begin{equation}\label{eq:explicitF}
 \mathcal F(u,v)=-(N-1)v^2+b(u)v+c(u),
\end{equation}
where
\begin{align}\label{eq:bc}
 b(u)&=(N+2)u^2+(-2jm-j-m+r)u+jm(1-r),\\
 c(u)&=u(j-u)(r+u)(m-u).\notag
\end{align}
For $0<u<j$, $c(u)>0$, so the positive stationary root is
\begin{equation}\label{eq:vplus}
 v_+(u)=\frac{b(u)+\sqrt{b(u)^2+4(N-1)c(u)}}{2(N-1)}.
\end{equation}
The Bernoulli admissibility constraint restricts this root to $v_+(u)\le u(j-u)/j$ at an actual stationary tilt.

For the identity in Lemma \ref{lem:identity}, substitute $u=(j+x)/2$, $B=j(L+1)/(4T)$, $T=N-1$, $L=N-j-1$, and $D=h-r$. Collecting powers of $D$ gives
\begin{align*}
 16T\mathcal F((j+x)/2,B)
 &=-AD^2+2x[A-j(L+1)]D-x^2[A+(L+1)E]\\
 &=-A(D-x)^2-(L+1)(2jDx+Ex^2),
\end{align*}
where $A=j^2(j-1)-Tx^2$ and $E=L(L+1)-2j$.

\subsection{The all-degree positivity polynomial}\label{app:poly}
The coefficient identity \eqref{eq:Kpositive} follows by setting $a=j-3$ and $L=a+2+S$. The three factors become
\[
 E=a(a+3)+(2a+5)S+S^2,\qquad
 (j-1)L-1=(a+1)(a+3)+(a+2)S,
\]
\[
 L(L+j)=(a+2+S)(2a+5+S).
\]
Therefore the polynomial is
\[
 [a(a+3)+(2a+5)S+S^2]^2[(a+1)(a+3)+(a+2)S]
 -4S^2(a+2+S)(2a+5+S).
\]
Multiplication and collection of $S^0,\ldots,S^5$ give exactly the six nonnegative coefficients displayed in \eqref{eq:Kpositive}.


\section{Proof of the Bernoulli optimal proxy}\label{app:bernoulli}
To prove \eqref{eq:bernproxy}, symmetry under $p\leftrightarrow1-p$ permits $0<p\le1/2$; the case $p=1/2$ is immediate from $\log\cosh(t/2)\le t^2/8$. Suppose $p<1/2$, write $q=1-p$ and $\alpha=\log(q/p)>0$, and set
\[
 f(t)=\log(q+pe^t)-pt,\qquad
 c=\frac{q-p}{2\alpha}.
\]
We claim $f(t)\le ct^2/2$ for all $t$, with equality at $t=0$ and $t=2\alpha$. Let
\[
 G(t)=\frac{ct^2}{2}-f(t).
\]
The identities
\begin{equation}\label{eq:bernref}
 G(2\alpha-t)=G(t),\qquad
 G'(0)=G'(\alpha)=0
\end{equation}
follow by elementary logarithm algebra and the fact that $f'(\alpha)=1/2-p=(q-p)/2$. Differentiating twice gives
\begin{equation}\label{eq:bernsecond}
 G''(t)=c-\frac1{4\cosh^2((t-\alpha)/2)}.
\end{equation}
The second term on the right is strictly increasing for $t<\alpha$ and strictly decreasing for $t>\alpha$. Furthermore $G''(0)>0$, since with $y=\alpha/2$,
\[
 4c=\frac{\tanh y}{y}>\operatorname{sech}^2 y.
\]
(The latter inequality follows by integrating the strictly decreasing function $\operatorname{sech}^2$ on $[0,y]$.) Thus $G''$ has one zero on $(0,\alpha)$, is positive initially, and negative near $\alpha$. As $G'(0)=G'(\alpha)=0$, it follows that $G'(t)>0$ on $(0,\alpha)$ and therefore $G(t)>0$ there. Reflection \eqref{eq:bernref} gives $G\ge0$ on $[\alpha,2\alpha]$ as well. For $t<0$ the variance term in \eqref{eq:bernsecond} is smaller than at zero, hence $G''(t)>0$; together with $G(0)=G'(0)=0$, this yields $G(t)>0$. Reflection handles $t>2\alpha$. Therefore $G(t)\ge0$ for all $t$, establishing the upper bound.

At $t=2\alpha$ equality holds, so $c$ is the least admissible variance proxy. Rewriting $c$ gives
\[
 \sigma^2_{\mathrm{opt}}(\Bern(p))
 =\frac{q-p}{2\log(q/p)}=
 \frac{1-2p}{2\log((1-p)/p)}.
\]
This proves \eqref{eq:bernproxy}. The limit as $p\to1/2$ is $1/4$ by l'H\^opital's rule.

\begin{thebibliography}{99}\setlength{\itemsep}{3pt}
\bibitem{ArbelMarchalNguyen} J. Arbel, O. Marchal and H. D. Nguyen, On strict sub-Gaussianity, optimal proxy variance and symmetry for bounded random variables, \emph{ESAIM: Probability and Statistics} \textbf{24} (2020), 39--55. \url{https://doi.org/10.1051/ps/2019018}.
\bibitem{AtouaniMarchalArbel} S. Atouani, O. Marchal and J. Arbel, Optimal sub-Gaussian variance proxy for 3-mass distributions, arXiv:2510.06132v1 (2025). \url{https://arxiv.org/abs/2510.06132}.
\bibitem{BardenetMaillard} R. Bardenet and O.-A. Maillard, Concentration inequalities for sampling without replacement, \emph{Bernoulli} \textbf{21} (2015), 1361--1385. \url{https://doi.org/10.3150/14-BEJ605}.
\bibitem{Barber} R. F. Barber, Hoeffding and Bernstein inequalities for weighted sums of exchangeable random variables, \emph{Electronic Communications in Probability} \textbf{29} (2024), 1--13. \url{https://doi.org/10.1214/24-ECP620}.
\bibitem{BarberErratum} R. F. Barber, Erratum to ``Hoeffding and Bernstein inequalities for weighted sums of exchangeable random variables'', \emph{Electronic Communications in Probability} \textbf{30} (2025), 1. \url{https://doi.org/10.1214/25-ECP715}.
\bibitem{BBL} J. Borcea, P. Br\"and\'en and T. M. Liggett, Negative dependence and the geometry of polynomials, \emph{Journal of the American Mathematical Society} \textbf{22} (2009), 521--567. \url{https://doi.org/10.1090/S0894-0347-08-00618-8}.
\bibitem{DiaconisFreedman} P. Diaconis and D. Freedman, Finite exchangeable sequences, \emph{Annals of Probability} \textbf{8} (1980), 745--764. \url{https://doi.org/10.1214/aop/1176994663}.
\bibitem{Fog} A. Fog, Sampling methods for Wallenius' and Fisher's noncentral hypergeometric distributions, \emph{Communications in Statistics---Simulation and Computation} \textbf{37} (2008), 241--257. \url{https://doi.org/10.1080/03610910701790236}.
\bibitem{GottschlingCaprio} N. M. Gottschling and M. Caprio, Hoeffding-Style Concentration Bounds for Exchangeable Random Variables, arXiv:2603.10190v2 (2026). \url{https://arxiv.org/abs/2603.10190v2}.
\bibitem{GreeneWellner} E. Greene and J. A. Wellner, Exponential bounds for the hypergeometric distribution, \emph{Bernoulli} \textbf{23} (2017), 1911--1950. \url{https://doi.org/10.3150/15-BEJ800}.
\bibitem{Hoeffding} W. Hoeffding, Probability inequalities for sums of bounded random variables, \emph{Journal of the American Statistical Association} \textbf{58} (1963), 13--30. \url{https://doi.org/10.1080/01621459.1963.10500830}.
\bibitem{Kou} S. G. Kou and Z. Ying, Asymptotics for a $2\times2$ table with fixed margins, \emph{Statistica Sinica} \textbf{6} (1996), 809--829. \url{https://www3.stat.sinica.edu.tw/statistica/oldpdf/A6n42.pdf}.
\bibitem{LeeKim} S. Lee and I. Kim, A sharper Hoeffding bound for weighted sums of exchangeable random variables, arXiv:2608.04900v2 (2026). \url{https://arxiv.org/abs/2608.04900v2}.
\bibitem{MarchalArbel} O. Marchal and J. Arbel, On the sub-Gaussianity of the Beta and Dirichlet distributions, \emph{Electronic Communications in Probability} \textbf{22} (2017), 1--14. \url{https://doi.org/10.1214/17-ECP92}.
\bibitem{Serfling} R. J. Serfling, Probability inequalities for the sum in sampling without replacement, \emph{Annals of Statistics} \textbf{2} (1974), 39--48. \url{https://doi.org/10.1214/aos/1176342611}.
\bibitem{Szego} G. Szeg\H{o}, Orthogonal Polynomials, \emph{American Mathematical Society Colloquium Publications}, vol. 23, fourth edition, AMS, 1975. \url{https://doi.org/10.1090/coll/023}.
\end{thebibliography}

\end{document}
